% !TEX program = pdflatex
% !TEX encoding = UTF-8 Unicode

% Plantilla de la clase `scrbook` del paquete KOMA-script para la
% elaboración de un TFG siguiendo las directrices del la comisión del
% Grado en Matemáticas de la Universidad de Granada.

% Francisco Torralbo Torralbo
% miércoles, 29 de abril de 2020

\documentclass{scrbook} 

\KOMAoptions{%
  fontsize=10pt,        % Tamaño de fuente
  paper=a4,             % Tamaño del papel
  headings=normal,      % Tamaño de letra para los títulos: small, normal, big
  % parskip=half,         % Espacio entre párrafos: full (una línea) o half (media línea)
  headsepline=false,    % Una linea separa la cabecera del texto
  cleardoublepage=empty,% No imprime cabecera ni pie en páginas en blanco 
  chapterprefix=false,  % No antepone el texto "capítulo" antes del número
  appendixprefix=false,	% No antepone el texto "Apéndice" antes de la letra
  listof=totoc,		    	% Añade a la tabla de contenidos la lista de tablas y figuras
  index=totoc,			    % Añade a la talba de contenidos una entrada para el índice
  bibliography=totoc,	  % Añade a la tabla de contenidos una entrada para bibliografía
  BCOR=5mm,           % Reserva margen interior para la encuadernación. 
                        % El valor dependerá el tipo de encuadernado y del grosor del libro.
  DIV=10,             % Cálcula el diseño de página según ciertos 
                        % parámetros. Al aumentar el número aumentamos el ancho de texto y disminuimos el ancho del margen. Una opción de 14 producirá márgenes estrechos y texto ancho.
}

% INFORMACIÓN PARA LA VERSIÓN IMPRESA
% Si el documento ha de ser impreso en papel de tamaño a4 pero el tamaño del documento (elegido en \KOMAoptions con la ocpión paper) no es a4 descomentar la línea que carga el paquete `crop` más abajo. El paquete crop se encargará de centrar el documento en un a4 e imprimir unas guías de corte. El procedimiento completo para imprenta sería el siguiente:
% 0. Determinar, según el tipo de encuadernación del documento, el ancho reservado para el proceso de encuadernación (preguntar en la imprenta), es decir, la anchura del área del papel que se pierde durante el proceso de encuadernación. Fijar la varibale BCOR de \KOMAoptions a dicho valor.
% 1. Descomentar la siguiente línea e imprimir una única página con las guías de corte
% 2. Cambiar la opción `cross` por `cam` (o `off`) en el paquete crop y volver a compilar. Imprimir el documento (las guías de corte impresas no inferfieren con el texto).
% 3. Usar la página con las guías impresas en el punto 1 para cortar todas las páginas.

% \usepackage[a4, odd, center, pdflatex, cross]{crop} % Permite imprimir el documento en un a4 (si el tamaño es más pequeño) mostrando unas guías de corte. Útil para imprenta.

% VERSIÓN ELECTRÓNICA PARA TABLETA
% Las opciones siguientes seleccionan un tamaño de impresión similar a una tableta de 9 pulgadas con márgenes estrechos. Útil para producir una versión en pdf para ser leída en una tableta en lugar de impresa.
% Para que la portada quede centrada correctamente hay que editar el
% archivo `portada.tex` y eliminar el entorno `addmargin`

% \KOMAoptions{fontsize=10pt, paper=19.7104cm:14.7828cm, twoside=false, BCOR=0cm, DIV=14}

% ---------------------------------------------------------------------
%	PAQUETES 
% ---------------------------------------------------------------------

% CODIFICACIÓN E IDIOMA
% ---------------------------------------------------------------------
\usepackage[utf8]{inputenc} 			    % Codificación de caracteres

% Selección del idioma: cargamos por defecto inglés y español (aunque este último es el idioma por defecto para el documento). Cuando queramos cambiar de idioma escribiremos:
% \selectlanguage{english} o \selectlanguage{spanish}

\usepackage[english, spanish, es-nodecimaldot, es-noindentfirst, es-tabla]{babel}

% Opciones cargadas para el paquete babel:
  % es-nodecimaldot: No cambia el punto decimal por una coma en modo matemático.
  % es-noindentfirst: No sangra los párrafos tras los títulos.
  % es-tabla: cambia el título del entorno `table` de "Cuadro" a "Tabla"

% Otras opciones del paquete spanish-babel:
  \unaccentedoperators % Desactiva los acentos en los operadores matemáticso (p.e. \lim, \max, ...). Eliminar esta opción si queremos que vayan acentuados

% MATEMÁTICAS
% ---------------------------------------------------------------------
\usepackage{amsmath, amsthm, amssymb} % Paquetes matemáticas
\usepackage{mathtools}                % Añade mejoras a amsmath
\mathtoolsset{showonlyrefs=true}      % sólo se numeran las ecuaciones que se usan
\usepackage[mathscr]{eucal} 					% Proporciona el comando \mathscr para
                                      % fuentes de tipo manuscrito en modo matemático sin sobreescribir el comando \mathcal

% TIPOGRAFÍA 
% ---------------------------------------------------------------------
% El paquete microtype mejora la tipografía del documento.
\usepackage[activate={true,nocompatibility},final,tracking=true,kerning=true,spacing=true,factor=1100,stretch=10,shrink=10]{microtype}

% Las tipografías elegidas para el documento, siguiendo la guía de estilo de la UGR,
% son las siguientes
% Normal font: 			URW Palladio typeface. 
% Sans-serif font: 	Gill Sans
% Monospace font: 	Inconsolata
\usepackage[T1]{fontenc}
\usepackage[sc, osf]{mathpazo} \linespread{1.05}         
\usepackage[scaled=.95,type1]{cabin} % sans serif in style of Gill Sans
% Si el paquete cabin da error usar el siguiente comando en su lugar
% \renewcommand{\sfdefault}{iwona} 
\usepackage{inconsolata}


% Selecciona el tipo de fuente para los títulos (capítulo, sección, subsección) del documento.
\setkomafont{disposition}{\sffamily\bfseries}

% Cambia el ancho de la cita. Al inicio de un capítulo podemos usar el comando \dictum[autor]{cita} para añadir una cita famosa de un autor.
\renewcommand{\dictumwidth}{0.45\textwidth} 

\recalctypearea % Necesario tras definir la tipografía a usar.

\usepackage{setspace}
% TABLAS, GRÁFICOS Y LISTADOS DE CÓDIGO
% ---------------------------------------------------------------------
\usepackage{booktabs}
% \renewcommand{\arraystretch}{1.5} % Aumenta el espacio vertical entre las filas de un entorno tabular

\usepackage{xcolor, graphicx}
% Carpeta donde buscar los archivos de imagen por defecto
\graphicspath{{img/}}

% IMAGEN DE LA PORTADA
% Existen varias opciones para la imagen de fondo de la portada del TFG. Todas ellas tienen en logotipo de la universidad de Granada en la cabecera. Las opciones son las siguientes:
% 1. portada-ugr y portada-ugr-color: diseño con marca de agua basada en el logo de la UGR (en escala de grises y color).
% 2. portada-ugr-sencilla y portada-ugr-sencilla-color: portada únicamente con el logotipo de la UGR en la cabecera.
\usepackage{eso-pic}
\newcommand\BackgroundPic{%
	\put(0,0){%
		\parbox[b][\paperheight]{\paperwidth}{%
			\vfill
			\centering
      % Indicar la imagen de fondo en el siguiente comando
			\includegraphics[width=\paperwidth,height=\paperheight,%
			keepaspectratio]{portada-ugr-sencilla-color}%
			\vfill
}}}

% \usepackage{listings} % Para la inclusión de trozos de código

\usepackage{longtable}

% CABECERAS
% ---------------------------------------------------------------------
% Si queremos modificar las cabeceras del documento podemos usar el paquete
% `scrlayer-scrpage` de KOMA-Script. Consultar la documentación al respecto.
% \usepackage[automark]{scrlayer-scrpage}

% VARIOS
% ---------------------------------------------------------------------

%\usepackage{showkeys}	% Muestra las etiquetas del documento. Útil para revisar las referencias cruzadas.

% ÍNDICE 
% Para generar el índice hay que compilar el documento con MakeIndex. Generalmente los editores se encargan de ello automáticamente.
% ----------------------------------------------------------------------
% \index{} para añadir un elemento
% \index{main!sub} para añadir un elementos "sub" bajo la categoría "main".
% \index{termino|textbf} para dar formato al número de página (negrita).
% \index{termino|see{termino relacionado}} para crear una referencia cruzada

% Ejemplo: \index{espacio homogéneo}, \index{superficie!mínima}, \index{esfera|see{espacio homogéneo}}

% Activar los siguientes comandos para generar el índice terminológico. Ver también comandos al final de este documento para incluir dicho índice en el pdf final.
% \usepackage{makeidx}
% \makeindex

% Para revisar las entradas al índice conforme las incluimos en el documento es útil el siguiente paquete. Conviene observar que mientras esté cargado no se generará el índice.
%\usepackage{showidx} % Muestra en el margen del documento las entradas añadidas al índice. Útil para revisar el documento. Si está activo el índice no se genera

% ---------------------------------------------------------------------
% COMANDOS Y ENTORNOS
% ---------------------------------------------------------------------
% Cargamos un archivo externo donde hemos incluido todos los comandos
% propios que vamos a usar en el documento.
% DEFINICIÓN DE COMANDOS Y ENTORNOS

% CONJUNTOS DE NÚMEROS

  \newcommand{\N}{\mathbb{N}}     % Naturales
  \newcommand{\R}{\mathbb{R}}     % Reales
  \newcommand{\Z}{\mathbb{Z}}     % Enteros
  \newcommand{\Q}{\mathbb{Q}}     % Racionales
  \newcommand{\C}{\mathbb{C}}     % Complejos

% TEOREMAS Y ENTORNOS ASOCIADOS

  % \newtheorem{theorem}{Theorem}[chapter]
  \newtheorem*{teorema*}{Teorema}
  \newtheorem{teorema}{Teorema}[chapter]
  \newtheorem{proposicion}{Proposición}[chapter]
  \newtheorem{lema}{Lema}[chapter]
  \newtheorem{corolario}{Corolario}[chapter]

    \theoremstyle{definition}
  \newtheorem{definicion}{Definición}[chapter]
  \newtheorem{ejemplo}{Ejemplo}[chapter]

    \theoremstyle{remark}
  \newtheorem{observacion}{Observación}[chapter]


% --------------------------------------------------------------------
% INFORMACIÓN DEL TFG Y EL AUTOR
% --------------------------------------------------------------------
\usepackage{xspace} % Para problemas de espaciado al definir comandos

\newcommand{\miTitulo}{El programa de Hilbert. Teoremas de Incompletitud de Gödel\xspace}
\newcommand{\miNombre}{David García Curbelo\xspace}
\newcommand{\miGrado}{Grado en Matemáticas}
\newcommand{\miFacultad}{Facultad de Ciencias}
\newcommand{\miUniversidad}{Universidad de Granada}
% Añadir tantos tutores como sea necesario separando cada uno de ellos
% mediante el comando `\medskip` y una línea en blanco
\newcommand{\miTutor}{
  Miguel Delgado Calvo-Flores \\ \emph{Ciencias de la Computación e Inteligencia Artificial} 

  % Añadir tantos tutores como sea necesario. 

}
\newcommand{\miCurso}{2021-2022\xspace}

% HYPERREFERENCES
% --------------------------------------------------------------------
\usepackage{xurl}
\usepackage[pagebackref]{hyperref}
% Opciones para el paquete hyperref
%----------------------------------

\hypersetup{%
  % hidelinks,            % Enlaces sin color ni borde. El borde no se imprime
  linkbordercolor=0.8 0 0,
  citebordercolor=0 0.8 0,
  citebordercolor=0 0.8 0,
  colorlinks = true,            % Color en texto de los enlaces. Comentar esta línea o cambiar `true` por `false` para imprimir el documento.
  linkcolor = [rgb]{0.5, 0, 0}, % Color de los enlaces internos
  urlcolor = [rgb]{0, 0, 0.5},  % Color de los hipervínculos
  citecolor = [rgb]{0, 0.5, 0}, % Color de las referencias bibliográficas
	pdftitle={\miTitulo},%
	pdfauthor={\textcopyright\ \miNombre, \miFacultad, \miUniversidad},%
  pdfsubject={Trabajo de fin de Grado},%
	pdfkeywords={},%
	pdfcreator={pdfLaTeX},%
}

% Redefinición del estilo para mostrar las referencias cruzadas en la bibliografía.
\renewcommand*{\backref}[1]{}
\renewcommand*{\backrefalt}[4]{{\footnotesize [%
    \ifcase #1 No citado%
    \or Citado en pág.~#2%
    \else Citado en págs. #2%
    \fi%
]}}

% Etiquetas en español para el comando \autoref
\def\chapterautorefname{Capítulo}
\def\appendixautorefname{Apéndice}
\def\sectionautorefname{Sección}
\def\subsectionautorefname{Subsección}
\def\figureautorefname{Fig.}
\def\tableautorefname{Tabla}

\def\teoremaautorefname{Teorema}
\def\proposicionautorefname{Proposición}
\def\lemaautorefname{Lema}
\def\corolarioautorefname{Corolario}
\def\definicionautorefname{Def.}
\def\observacionautorefname{Observación}
\def\ejemploautorefname{E.j.}

% Pone automáticamente un parántesis para las ecuaciones
\def\equationautorefname~#1\null{Ec.~(#1)\null}


\begin{document}

% --------------------------------------------------------------------
% FRONTMATTER
% -------------------------------------------------------------------
\frontmatter % Desactiva la numeración de capítulos y usa numeración romana para las páginas

% \pagestyle{plain} % No imprime cabeceras

% !TeX root = ../libro.tex
% !TeX encoding = utf8

%*******************************************************
% Titlepage
%*******************************************************
\begin{titlepage}
  \AddToShipoutPicture*{\BackgroundPic}
  \phantomsection 
  \pdfbookmark[1]{Título}{title}

  % Para que el título esté centrado en la página.
  % Los valores numéricos deberán elegirse de acuerdo con el diseño de
  % página (sobre todo si se cambia la opción BCOR o DIV).
  \begin{addmargin}[2.575cm]{0cm}
  \begin{flushleft}
    \Large  
    \hfill\vfil

    \large{\miFacultad}
    \vfill

    {\large\textsc\miGrado} \vfill


    {\large\textsc{trabajo de fin de grado}}

    \begin{flushleft}
      \Huge
      \setstretch{0.8}
      \miTitulo
    \end{flushleft}

    \vfill\vfill\vfill\vfill

    \textsf{\normalsize{Presentado por:}}\\
    {\normalsize\textrm{\miNombre}} 
    \bigskip

    \textsf{\normalsize{Tutor:}}\\
    {\normalsize\rmfamily\miTutor}

    \bigskip
    \textsf{\normalsize{Curso académico \miCurso}}
  \end{flushleft}  
  \end{addmargin}       

\end{titlepage}   
\cleardoublepage
\endinput
                    
% !TeX root = ../libro.tex
% !TeX encoding = utf8

%*******************************************************
% Little Dirty Titlepage
%*******************************************************

\thispagestyle{empty}

\begin{center}
  \large  

  \vspace*{\stretch{1}}

  \begingroup
  \huge{\miTitulo} \\ \bigskip
  \endgroup

  \textrm{\miNombre}

  \vspace{\stretch{5}}

\end{center}  

\newpage
\thispagestyle{empty}

\hfill

\vfill

\miNombre \textit{\miTitulo}.

Trabajo de fin de Grado. Curso académico \miCurso.
\bigskip

\begin{minipage}[t]{0.25\textwidth}
  \flushleft
  \textbf{Responsable de tutorización}
\end{minipage}
\begin{minipage}[t]{0.45\textwidth}
  \flushleft
  \miTutor
\end{minipage}
\begin{minipage}[t]{0.30\textwidth}
  \flushright
  \miGrado
  \medskip

  \miFacultad
  \medskip

  \miUniversidad
\end{minipage}

\newpage
\endinput
                     
% !TeX root = ../libro.tex
% !TeX encoding = utf8
%
%*******************************************************
% Declaración de originalidad
%*******************************************************

\thispagestyle{empty}

\hfill\vfill

\textsc{Declaración de originalidad}\\\bigskip

D./Dña. \miNombre \\\medskip

Declaro explícitamente que el trabajo presentado como Trabajo de Fin de Grado (TFG), correspondiente al curso académico \miCurso, es original, entendida esta, en el sentido de que no ha utilizado para la elaboración del trabajo fuentes sin citarlas debidamente.
\medskip

En Granada a \today 
\begin{flushleft} 
Fdo: \miNombre 

\end{flushleft}

\vfill

\cleardoublepage
\endinput
   
% !TeX root = ../libro.tex
% !TeX encoding = utf8

%*******************************************************
% Dedication
%*******************************************************
\thispagestyle{empty}
\phantomsection 
\pdfbookmark[1]{Dedicatoria}{Dedicatoria}

\hfill
\vfill

\begin{flushright}
\itshape
A mi buen amigo Diego, \\
por haberme hecho ver las matemáticas 
no son sólo cosa de locos.
\end{flushright}

\vfill

\cleardoublepage
\endinput
                % Opcional
% !TeX root = ../libro.tex
% !TeX encoding = utf8

%*******************************************************
% Table of Contents
%*******************************************************
\phantomsection
\pdfbookmark[0]{\contentsname}{toc}

\setcounter{tocdepth}{2} % <-- 2 includes up to subsections in the ToC
\setcounter{secnumdepth}{3} % <-- 3 numbers up to subsubsections

% \manualmark
% \markboth{\textsc{\contentsname}}{\textsc{\contentsname}}
\tableofcontents 

%*******************************************************
% List of Figures and of the Tables
%*******************************************************

    % *******************************************************
    %  List of Figures
    % *******************************************************    
    \phantomsection 
    % \listoffigures

    %*******************************************************
    % List of Tables
    %*******************************************************
    \phantomsection 
    % \listoftables
    
    %*******************************************************
    % List of Listings
    % The package \usepackage{listings} is needed
    %*******************************************************      
	  % \phantomsection 
    % \renewcommand{\lstlistlistingname}{Listados de código}
    % \lstlistoflistings 

\cleardoublepage
            
% % !TeX root = ../libro.tex
% !TeX encoding = utf8

%*******************************************************
% Agradecimientos
%*******************************************************

\chapter{Agradecimientos}

Agradecimientos del libro (opcional, ver archivo \texttt{preliminares/agradecimiento.tex}).

\cleardoublepage
\endinput
            % Opcional

% \pagestyle{scrheadings} % A partir de ahora sí imprime cabeceras

% !TeX root = ../libro.tex
% !TeX encoding = utf8
%
%*******************************************************
% Summary
%*******************************************************

\selectlanguage{english}
\chapter{Summary}

This paper presents the history of the theorem that shook the foundations of all known mathematics: Gödel's incompleteness theorems. 
To do so, we will give a brief historical review of what motivated this great mathematician of the twentieth century to deal with this 
peculiar subject, in which we will see that this merit is attributed to David Hilbert and his list of 23 problems. 

At the end of the 19th century, mathematics felt a slight tremor in its foundations, giving rise to the foundational crisis of mathematics, 
and it is here that Hilbert appears with his second problem of the list, in which he poses as an urgent problem the axiomatization of arithmetic. 
What Hilbert did not know was that this problem would never be solved. 

Gödel, after graduating in mathematics, comes across a recently published book by Hilbert and Ackermann in which an open problem of logic is posed: 
proving the completeness of the first-order logical calculus. Emerging from the unknown world, this great up-and-coming mathematician proposes 
an irrefutable solution in less than a year, and in a paper of barely 10 pages in which he states his strong result: 

\begin{center}
    \textit{In a first-order logic, every formula that is valid in a logical sense is provable.}
\end{center}
 
This simple statement promises us the existence of demonstration of any statement of elementary logical formulas, provided that these are in 
infinite numerable quantity. This great result gave much hope to Hilbert, who had hitherto held high hopes for the possibility of solving his 
second problem proposed at the 1900 congress. 

However, what Hilbert could not foresee was that, a year later, Gödel would publish the result that would shake the whole of mathematics, 
and not only that of his time, but forever. In one of the most important articles of twentieth century mathematics, Gödel presents the well-known 
Incompleteness Theorems, by means of a constructive proof based on a gödelization, which is nothing more than an assignment of numbers to 
elements of first-order logic. In them he states as follows:

\begin{center}
    \begin{enumerate}
        \item \textit{No formal mathematical theory is consistent and complete.}
        \item \textit{If the system of axioms of such a formal system is consistent, such consistency cannot be proved by means of these axioms.}
    \end{enumerate}
\end{center}

This result represents a great wall in the advancement of the knowledge of mathematics. In the first instance, we will never be able to create a 
mathematical system in which we can be sure of being able to prove every result, and also that we will not have contradictions in it. 

In this small bibliographic work we intend to collect the original proofs published between 1929 and 1931 by Kurt Gödel himself 
(although later other simpler and improved proposals have been published, I believe that in the original ones we can better understand the 
objectives of the same through a constructive process), in addition to the historical context necessary to understand the motivation of these 
proofs.

% Al finalizar el resumen en inglés, volvemos a seleccionar el idioma español para el documento
\selectlanguage{spanish} 
\endinput
                    
% !TeX root = ../libro.tex
% !TeX encoding = utf8
%
%*******************************************************
% Introducción
%*******************************************************

% \manualmark
% \markboth{\textsc{Introducción}}{\textsc{Introducción}} 

\chapter{Introducción}

En este trabajo se presenta la historia del teorema que hizo tambalear los cimientos de toda la matemática conocida:
los teoremas de incompletitud de Gödel. Para ello daremos un pequeño repaso histórico acerca de qué fue lo 
que motivó a este gran matemático del siglo XX a tratar este tema tan peculiar, en el que veremos que dicho 
mérito se le atribuye a David Hilbert y a su lista de 23 problemas.

A finales de siglo XIX, la matemática  siente un pequeño temblor en sus fundamentos, dando lugar a la crisis fundacional
de la matemática, y es aquí cuando aparece Hilbert con su segundo problema de la lista, en el cual plantea como 
problema urgente la axiomatización de la aritmética. Lo que Hilbert desconocía es que este problema nunca tendría solución.

Gödel, después de haberse licenciado en matemáticas, se encuentra con un libro recién publicado de Hilbert y Ackermann en el que se plantea
un problema abierto de la lógica: demostrar la completitud del cálculo lógico 
de primer orden. Emergiendo del mundo desconocido, este gran matemático prometedor propone una solución irrefutable en menos de un año, 
y en un artículo de apenas 10 páginas en el que enuncia su resultado fuerte:

\begin{center}
    \textit{En una lógica de primer orden, toda fórmula que es válida en un sentido lógico es demostrable.}
\end{center}

Este enunciado tan simple nos promete la existencia de demostración de cualquier enunciado de fórmulas lógicas elementales, siempre
que estas se encuentren en cantidad infinita numerable. Este gran resultado dio muchas esperanzas a Hilbert, el cual sostenía 
hasta entonces grandes esperanzas acerca de la posibilidad de resolver su segundo problema propuesto en el congreso de 1900.

Sin embargo, lo que Hilbert no pudo prevenir es que, un año después, Gödel publicaría el resultado que haría temblar a toda la matemática, 
y no sólo la de su época, sino para siempre. En uno de los artículos más importantes de la matemática del siglo XX, Gödel presenta los 
conocidos \textit{Teoremas de incompletitud}, mediante una prueba constructiva basada en una \textit{gödelización}, que no es más que una asignación
de números a elementos de la lógica de primer orden. En ellos enuncia como sigue:

\begin{center}
    \begin{enumerate}
        \item \textit{Ninguna teoría matemática formal es consistente y completa}
        \item \textit{Si el sistema de axiomas de dicho sistema formal es consistente, dicha consistencia no se puede probar mediante dichos axiomas.}
    \end{enumerate}
\end{center}

Este resultado supone un gran muro en el avance del conocimiento de la matemática. En una primera instancia, nunca podremos de ninguna manera
crear un sistema matemático en el que tengamos la certeza de poder probar todo resultado, y además tampoco de que no vayamos a tener contradicciones en el mismo.\\

En este pequeño trabajo bibliográfico se pretende recoger las demostraciones originales publicadas entre los años 1929 y 1931 por el propio Kurt Gödel (aunque 
posteriormente se hayan publicado otras propuestas más simples y mejoradas, considero que en las originales se pueden entender mejor los objetivos de la misma mediante
un proceso constructivo), además del contexto histórico necesario para entender la motivación de dichas demostraciones.


\endinput
               

% --------------------------------------------------------------------
% MAINMATTER
% --------------------------------------------------------------------
\mainmatter % activa la numeración de capítulos, resetea la numeración de las páginas y usa números arábigos

\setpartpreamble[c][0.75\linewidth]{%
	\bigskip % Deja un espacio vertical en la parte superior
  % Aqui metemos resumen del programa de hilbert
}
\part{El programa de Hilbert} % Dividir un libro en partes OPCIONAL

% !TeX root = ../libro.tex
% !TeX encoding = utf8

\setchapterpreamble[c][0.75\linewidth]{% 
	\sffamily
  % Al inicio de cada capítulo puede incluirse un breve resumen. Esto es opcional.
	\dictum[David Hilbert]{¿Quién de nosotros no se alegraría de levantar el velo tras el que se oculta el futuro; de echar una mirada a los próximos avances de nuestra ciencia y a los secretos de su desarrollo durante los siglos futuros?}

  \par\bigskip
}
\chapter{Una cortina hacia el futuro}\label{ch:Hilbert-1}

Con estas ambiciosas palabras comenzaba David Hilbert su conferencia desde el estrado de la universidad de La Sorbona, París, dando
así comienzo al Congreso Internacional de Matemáticos de 1900.
Dicha conferencia se distinguiría notablemente del resto de Congresos en los que ya había participado, debido a su gran influencia 
en la matemática de la época y a la ambición de la misma, que ha conseguido dejarnos retales hasta la actualidad. 


\section{Los objetivos}

Hilbert tuvo como principal reto en su discurso el crear un punto de inflexión en la matemática de la época, tratando de trazar 
una dirección y un sentido hacia el avance y el futuro de la matemática. Dicho objetivo lo abordó mediante la propuesta de una lista de 
23 problemas, ya que el propio Hilbert afirmaba que las matemáticas avanzaban proponiendo problemas, y que éstos en sí mismos 
son un signo de que la disciplina sigue viva.

Dicho programa ha tenido tal sonoridad y repercusión en la matemática del siglo XX debido tanto a la diversidad de campos que abarca, 
como a la propia influencia del conferenciante. Tal ha sido dicha importancia, que varios de estos problemas están aún sin resolver, 
en vistas de ideas y planteamientos matemáticos que aún no han sido organizados de manera correcta, o inclusive de nuevos campos y 
descubrimientos ocultos tras el velo de la ignorancia.


\section{La lista}

La recopilación de dichos 23 problemas abarca no sólo el campo que nos concierne, sino también otros tangenciales que eran de especial 
interés para Hilbert, como lo era la física. En este campo presentó especial interés gracias en gran medida a su buen amigo Hermann Minkowski,
y gracias a él estuvo activo incluso muchos años después de la muerte de su compañero, participando en los principales avances de la física del momento.

Dentro de los problemas de las matemáticas que presenta, podemos ver varios campos de especial interés en la matemática de la época, tocando geometrías no euclídeas
(todo un campo revolucionario de la época), un comienzo de atisbo en el análisis funcional (fundado en gran parte por el propio David Hilbert),
o incluso problemas que afrontan la crisis fundacional de la matemática, un lastre del que la matemática parecía no poder librarse, pero en el que Hilbert 
tenía mucha fe y seguridad.

Como comentábamos anteriormente, Hilbert era uno de los grandes matemáticos de la época, y debido a ésto, en sus estudios trabajaba con la 
matemática latente de la época. Por ello, en la lista de problemas podemos observar que gran número de ellos (alrededor de tres cuartas partes del
programa) están constituidos por problemas de campos que el propio Hilbert trataba, y que por tanto, debían ser temas procedentes del motor de avance
y del posible futuro que pudiera tomar la matemática debido al programa. 

  
\subsection{Los 23 problemas}

En la siguiente \autoref{tb:programa} se presenta una lista de los 23 problemas anunciados por Hilbert en el Congreso Internacional de Matemáticos en París,
aunque dicha lista no fuera publicada hasta un tiempo después. Como se puede ver, muchos de ellos carecen de una explicación clara y concisa del objetivo del problema,
y por éste motivo se ha considerado que muchos de ellos eran para Hilbert una dirección que tomar en el avance de la matemática, más que un problema como tal.

\begin{longtable}{|c|p{7cm}|c|} \toprule
  \centering
  % \begin{tabular}{|c|p{7cm}|c|} 
    Problema & Enunciado & Estado \\ \midrule \midrule
    1 & La hipótesis del continuo & Solución parcial \\ \midrule 
    2 & La compatibilidad de los axiomas de la aritmética & Resuelto \\ \midrule 
    3 & La igualdad de los volúmenes de dos 
        tetraedros de la misma base y la misma altura & Resuelto \\ \midrule 
    4 & Construcción de todas las métricas cuyas rectas sean geodésicas & Demasiado genérico \\ \midrule
    5 & El concepto de Lie de grupo continuo de transformaciones sin 
        la hipótesis de la diferenciabilidad de las funciones que 
        definen el grupo & Resuelto \\ \midrule
    6 & Tratamiento matemático de los axiomas de la física & Solución parcial \\ \midrule
    7 & Irracionalidad y trascendencia de ciertos números & Resuelto \\ \midrule
    8 & Problemas de números primos & Sin resolver \\ \midrule
    9 & Demostración de la ley de reciprocidad más 
        general en cualquier campo de números & Parcialmente resuelto \\ \midrule
    10 & Determinación de la resolubilidadde la ecuación diofántica & Resuelto \\ \midrule
    11 & Formas cuadráticas con coeficientes numéricos 
        algebraicos cualesquiera & Parcialmente resuelto \\ \midrule
    12 & Extensión del teorema de Kronecker sobre campos abelianos 
        a cualquier dominio de racionalidad algebraico & Sin resolver \\ \midrule
    13 & Imposibilidad de la solución general de 7º
        grado por medio de funciones de sólo dos argumentos & Resuelto \\ \midrule
    14 & Demostración de la finitud de ciertos sistemas completos de funciones & Resuelto \\ \midrule
    15 & Fundamentación rigurosa del cálculo enumerativo de Schubert & Solución parcial\\ \midrule
    16 & Problema de la topología de curvas y superficies algebraicas & Sin resolver \\ \midrule
    17 & Expresión de formas definidas por cuadrados & Resuelto \\ \midrule
    18 & Construcción del espacio a partir de poliedros congruentes & Resuelto \\ \midrule
    19 & ¿Son siempre necesariamente analíticas las soluciones de 
        problemas regulares en el cálculo de variaciones? & Resuelto \\ \midrule
    20 & El problema general de los valores de contorno & Resuelto \\ \midrule
    21 & Demostración de la existencia de ecuaciones diferenciales
        lineales que tienen prescrito un grupo monodrómico & Resuelto \\ \midrule
    22 & Uniformización de relaciones analíticas 
        por medio de funciones automorfas & Resuelto \\ \midrule
    23 & Desarrollo adicional de los métodos del cálculo de variaciones & Sin resolver \\ \bottomrule
  % \end{tabular}
  \caption{Lista de los 23 problemas del Programa}
  \label{tb:programa}
\end{longtable}

En la tabla podemos ver que muchos de los temas que trató no son un enunciado a un problema como tal, sino más bien a un campo de investigación, el cual en su 
discurso recogía varios problemas a resolver o simplemente unas directrices que seguir.



\section{El escenario previo}

No todos los problemas fueron motivo de estudio de Hilbert. Algunos, como los problemas que tienen que ver con la función zeta de Riemann o el teorema de uniformización,
ya eran problemas resonantes de la época que no se podían ignorar como problemas relevantes para el próximo siglo. En éste sentido Hilbert quiso hacer un compendio, 
no sólo de problemas de su propio campo e inquietudes, sino también de las grandes obviedades de la época. Por éste motivo uno también puede intuir otros posibles problemas
que hubieran podido ser motivo de estudio, y que sin embargo Hilbert ignoró\footnote{Un ejemplo resonante es el problema de dar una definición formal de integral de una función.}.

Pero entre los problemas que ataviaban a la sociedad matemática de la época, uno de los temas más recurrentes era la crisis fundacional. Nos encotramos ante un marco histórico 
en el que resultados recientes hacen temblar los cimientos de la matemática\footnote{Cabe destacar como ejemplos resonantes a la teoría de conjuntos de Cantor y a la paradoja 
de Russel, entre otros.}, y por consiguiente todo campo natural sedimentado en ella. Es por ello por lo que Hilbert reclama parte del protagonismo en la escena y propone dos problemas
de fundamentación axiomática: la axiomatización de la aritmética y de la física\footnote{Hilbert tenía gran confianza y certeza en que la axiomatización de la aritmética y la demostración de 
completitud de la misma era cuestión de tiempo, y por tanto también propuso la axiomatización de la física, que era el campo en el que también estaba enfocado. Véase \cite{kreisel1976have}.}.
Éste ve como necesario e imprescindible la resolución de dicho problema en el que, a pesar de tener gran certeza en la veracidad de su consistencia, era necesario una demostración rigurosa
y completa, que evitara la futura incertidumbre de posibles contradicciones no resolubles, y que quitaran de todo sentido a las matemáticas\footnote{En este sentido se puede ver más profundamente 
en palabras de Morris Kline, en \cite{kline2000matematicas}.}.

Es por éste problema (que trataremos más adelante) junto con muchos otros, que fueron motivación para Hilbert para incluirlos en su programa, y no sólo ganar más resonancia con ellos sino 
también tratar de motivar en su solución para que el futuro de la matemática pisara con paso firme, como toda la sociedad matemática deseaba en aquella época.

\chapter{Una mancha en la cortina} \label{ch:Hilbert-2}

\section{Los defectos del programa}

Debido a la gran influencia de David Hilbert, el programa era algo que no se podía refutar ni sacar del plano principal de estudio para las próximas generaciones. Sin embargo,
no por ello significa que el programa no tuviera objeciones. En efecto, muchos de los "problemas" propuestos no fueron más que un disparo al aire; una forma de que la matemática 
de la época mirara hacia donde Hilbert quería mirar. 

Entre los problemas planteados se encuentran tanto problemas del propio estudio de Hilbert, como problemas relevantes de la época, e incluso incluyó intereses que él quería que se
investigasen, pero que desgraciadamente no pertenecían a su campo de estudio. Ésto llevó a que quedaran problemas enunciados de forma demasiado vaga, lo que provocó un desinterés 
general en la sociedad y que, por ello, no fueran influyentes en el desarrollo posterior de la materia. Esto fue lo que ocurrió con el bloque de problemas del 13 al 18, siendo dentro
de los 23 los menos coherentes, y sobre los que Hilbert tenía el tacto menos seguro.\footnote{Esto sin embargo no ocurrió con los últimos 5 problemas, en los que Hilbert tenía terreno
más firme, e incluso había avanzado en algunos de ellos. Véase ~\cite[Pág. 85]{BREZIS199876}.}  


\section{La axiomatización DE LOURDES (mi segunda jefa)} MAMA NO QUIERO ESTUDIAR MAS

El problema de la axiomatización de la aritmética (enunciado en el segundo problema), no tuvo la respuesta esperada que Hilbert quería. En efecto, la axiomatización de la aritmética 
era algo estrictamente necesario, pero el escepticismo de la sociedad de la época ponía en duda la mera posibilidad de cualquier avance posible.

Cierto es que en los años siguientes se dieron pequeños pasos en el avance de la consistencia de la aritmética, tales como la obra de Russel y Whitehead con los \textit{Principia Mathematica}
y su teoría de tipos\footnote{Véase \cite{an1910principia}.}, o la emergente teoría de conjuntos de Cantor. Se pudo ver en ambas que, mediante cada una de ellas por separado, se podía probar 
que ambos sistemas formales podían expresar toda la matemática conocida.



Tambien puedo hablar aqui de los avances de hilbert con el libro que se habla en la intro de Godel.


\section{La llegada de Gödel a escena} 
\dictum[Samuel Beckett]{Sam enunció: "Lo que va a decir es falso" \\ Kurt replicó: "Lo que acaba de decir es verdadero"}

% Aquí introducción histórica y formal del cálculo lógico de primer orden
En verano de 1928, el joven y recién licenciado en matemáticas Kurt Gödel, se encuentra con un libro que le hace entrar en escena: \textit{Elementos de la lógica teórica}, de 
Hilbert y Ackermann\footnote{Publicado en 1928, siendo el primer libro de texto de lógica en el sentido actual. Véase \cite{hilbert1962elementos}.}. Dicho libro no sólo se presenta un cálculo deductivo, 
sino que por primera vez se plantea la cuestión de si ese cálculo es completo o no, dejando dicho problema abierto. Debido a la inquietud que causó en la matemática de la época,
y del impacto y la novedad que causó en el propio Gödel, fue motivo suficiente para tomar como tema para su tesis doctoral el tema de la completitud del cálculo lógico de primer orden.

Para Gödel, el campo de estudio, a pesar de ser de su interés, era completamente nuevo para él, el cual le planteaba como reto a largo plazo la resolución del segundo problema
de Hilbert en el intento de formalización de la aritmética. 

Como vimos, Hilbert ya había hecho pequeños avances en el tema, llegando a construir un cálculo deductivo de primer orden, el cual fue celebrado por toda la sociedad matemática 
como un gran avance en la solución del problema. 


%\section{Bibliografía}

%\cite{Zach_2007,Fer_Castro_2020,gray2000hilbert,stewart1987problems,gray2003reto,alma991014322964704990}



\endinput

% Añadir tantos capítulos como sea necesario

\setpartpreamble[c][0.75\linewidth]{%
	\bigskip % Deja un espacio vertical en la parte superior
  Aqui metemos resumen de los teoremas
}

\part{Los teoremas de incompletitud de Gödel}

% !TeX root = ../libro.tex
% !TeX encoding = utf8

\chapter{La completitud del cálculo lógico de primer orden}

\section{Definiciones y lemas previos}

Definamos previamente el sistema axiomático sobre el que vamos a trabajar:
\footnote{Coincide, exceptuando el principio de asociatividad (que es redundante), con lo expuesto en los apartados 1 y 10 de 
\textit{Principia Mathematica} (Véase \cite{an1910principia})}
\begin{itemize}
    \item Signos básicos primitivos:\footnote{A partir de ellos pueden definirse $\wedge$, $\rightarrow$, $\leftrightarrow$ y 
            $\exists$ del modo habitual.}
            \begin{itemize}
                \item $\lnot$
                \item $\vee$
                \item $\forall$
            \end{itemize}
    \item Axiomas formales:
            \begin{enumerate}
                \item $ \text{X} \vee \text{X} \rightarrow \text{X} $
                \item $ \text{X} \rightarrow \text{X} \vee \text{Y} $
                \item $ \text{Y} \vee \text{X} \rightarrow \text{X} \vee \text{Y} $
                \item $ (\text{X} \rightarrow \text{Y}) \rightarrow (\text{Z} \vee \text{X} \rightarrow \text{Z} \vee \text{Y}) $
                \item $ \forall \text{x Px} \rightarrow \text{Py} $
                \item $ \forall \text{x } (\text{X} \vee \text{Px}) \rightarrow \text{X} \vee \forall \text{x Px} $
            \end{enumerate}
    \item Reglas de inferencia: \footnote{No todas ellas están explícitamente formuladas en los resultados de Russel y Whitehead,
            pero todas ellas son usadas continuamente en sus deducciones.}
            \begin{enumerate}
                \item El esquema de inferencia: De $\alpha$ y $\alpha \rightarrow \beta$ se puede inferir $\beta$.
                \item La regla de sustitución para variables sentenciales y predicativas.
                \item De $\alpha(\text{x})$ puede inferirse $\forall \text{x } \alpha(\text{x})$.
                \item Unas variables individuales (libres o ligadas) pueden ser reemplazadas por cualesquiera otras, con tal de que 
                        ello no produzca ningún solapamiento de los alcances de las variables designadas mediante el mismo signo.
            \end{enumerate}
\end{itemize}

Para las deducciones que van a proceder, también es conveniente establecer algunas abreviaturas en calidad de notación:

\begin{enumerate}
    \item Las notaciones $\pi_1, \thinspace \pi_2, \thinspace \pi_3, \thinspace \rho, $ ... designan prefijos\footnote{Con prefijos 
    se hace referencia a conjuntos de signos básicos primitivos.} cualesquiera, es decir, filas de signos
            finitas de la forma $\forall \text{x} \exists \text{y}, \thinspace \forall \text{x} \forall \text{y} \exists \text{z} \forall \text{u},$ ...
    \item Las letras alemanas minúsculas $\mathfrak{x, \thinspace y, \thinspace u, \thinspace v,} $ etc., designan \textit{n-tuplos} de variables
            individuales, es decir, filas de signos del tipo $\text{xyz}, \thinspace \text{x}_2\text{x}_1\text{x}_2\text{x}_3,$ ..., 
            donde la misma variable puede aparecer varias veces en la misma sentencia.
            \begin{itemize}
                \item[Nota:] Del modo correspondiente hay que entender sentencias como $\forall \mathfrak{x} \thinspace \exists \mathfrak{y},$ etc. 
                            Si una misma variable aparece varias veces en $\mathfrak{x}$, hay que pensar, naturamente, que sólo está escrita una vez
                            en $\forall \mathfrak{x} \thinspace \exists \mathfrak{y},$ ...
            \end{itemize}
\end{enumerate}

Consideremos a continuación una serie de lemas, para las que no presentamos demostración por no ser relevantes para el desarrollo del resto de resultados 
(además de ser demostraciones relativamente sencillas de realizar).

\begin{lema} \label{lem:lema-1}
    Para cada n-tuplo $\mathfrak{x}$ los siguientes resultados son deducibles:
    \begin{enumerate}
        \item[$a)$] $ \forall \mathfrak{x} \text{F} \mathfrak{x} \rightarrow \exists \mathfrak{x} \text{F} \mathfrak{x}  $
        \item[$b)$] $ \forall \mathfrak{x} \text{F} \mathfrak{x} \wedge \exists \mathfrak{x} \text{G} \mathfrak{x} \rightarrow \exists \mathfrak{x} 
                    (\text{F} \mathfrak{x} \wedge \text{G} \mathfrak{x}) $
        \item[$c)$] $ \forall \mathfrak{x} \lnot \text{F} \mathfrak{x} \leftrightarrow \exists \mathfrak{x} \text{F} \mathfrak{x} $
    \end{enumerate}
\end{lema}

\begin{lema} \label{lem:lema-2}
    Si $\mathfrak{x}$ y $\mathfrak{x}'$ sólo se diferencian por el orden en que están escritas las variables, entonces el siguente enunciado es deducible:
    $$\exists \mathfrak{x} \text{F} \mathfrak{x} \rightarrow \exists \mathfrak{x}' \text{F} \mathfrak{x}'$$ 
\end{lema}

\begin{lema} \label{lem:lema-3}
    Si todas las variables $\mathfrak{x}$ son distintas entre sí y $\mathfrak{x}'$ tiene el mismo número de miembros que $\mathfrak{x}$, entonces el siguiente enunciado
    es deducible:
    $$\forall \mathfrak{x} \text{F} \mathfrak{x} \rightarrow \forall \mathfrak{x}' \text{F} \mathfrak{x}'$$
\end{lema}

\begin{lema} \label{lem:lema-4}
    Si $\pi_i$ designa uno de los prefijos $\forall x_i$, $\exists x_i$; y $\rho_i$ designa uno de los prefijos $\forall y_i$, $\exists y_i$, 
    entonces el siguiente es deducible:\footnote{Un resultado análogo vale para $\vee$ en vez de para $\wedge$.}
    $$ \pi_1 \pi_2 \cdots \pi_n Fx_1 x_2 \cdots x_n \wedge \rho_1 \rho_2 \cdots \rho_m G y_1 y_2 \cdots y_m \leftrightarrow \pi(Fx_1 x_2 \cdots x_n \wedge Gy_1 y_2 \cdots y_m ) $$
    para cada prefijo $\pi$ que se componga de los $\pi_i$ y $\rho_i$ y que satisfaga la condición de que $\pi_i$ esté delante de $\pi_k$ 
    para $i < k \leq n$ y de que $\rho_i$ esté delante de $\rho_k$ para $i < k \leq m$ 
\end{lema}

\begin{lema} \label{lem:lema-5}
    Toda fórmula puede ponerse en forma normal prenexa, es decir, para cada fórmula $\alpha$ hay una fórmula prenexa $\gamma$, tal que $\alpha \leftrightarrow \gamma$ es deducible.
\end{lema}

\begin{lema} \label{lem:lema-6}
    Si $\alpha \leftrightarrow \beta$ es deducible, entonces también lo es $\varphi(\alpha) \leftrightarrow \varphi(\beta)$, 
    donde $\varphi(\alpha)$ designa una fórmula cualquiera que contenga $\alpha$ como parte.\footnote{Estos dos últimos resultados se pueden estudiar 
    en detalle en la tercera sección de \cite{hilbert1962elementos}.}  
\end{lema}

\begin{lema} \label{lem:lema-7}
    Cada fórmula conectiva válida es deducible, es decir, los axiomas 1-4 constituyen un sistema suficiente de axiomas para el cálculo conectivo.
\end{lema}

Con estos resultados previos ya estamos en condiciones de afrontar el problema que nos concierne.


%%%%%%%%%%%%%%%%%%%%%%%%%%%%%%%%%%%%%
\section{Exposición y demostraciones}

El teorema de completitud semántica de la lógica de primer orden aparece en el artículo de Gödel de 1930 como Teorema I:

\begin{teorema}\label{thm:teoremaI}
    Cada fórmula válida de la lógica de primer orden es deducible.
\end{teorema}

El presente teorema, objeto principal de estudio de esta sección, sería trivialmente demostrable si pudiesemos probar
el siguiente:
\begin{teorema}\label{thm:teoremaII}
    Cada fórmula de la lógica de primer orden es o refutable
    \footnote{<<$\varphi$ es refutable>> significa <<$\lnot \varphi$ es deducible>>.} o satisfacible (sobre un universo infinito numerable).
\end{teorema}

Y por ello surge el siguiente resultado: 

\begin{proposicion}
    \autoref{thm:teoremaII} $\Rightarrow$ \autoref{thm:teoremaI}
\end{proposicion}
\begin{proof}
    Sea $\alpha$ una fórmula válida. Siendo esto así, $\lnot \alpha$ no es satisfacible, y aplicando \autoref{thm:teoremaII} 
    tenemos que $\alpha$ es refutable. Por tanto, con ello se tiene que $\lnot \lnot \alpha$ (y como consecuencia, también $\alpha$) es una fórmula
    deducible.\footnote{El recíproco del anterior resultado también es cierto y con una demostración igual de simple, aunque no la demostraremos por no ser relevante
    en la demostración del \autoref{thm:teoremaI}.}
\end{proof}

\begin{definicion}
    Una \textit{K-fórmula} es una fórmula $\kappa$ perteneciente a una clase del conjunto de fórmulas $K$ cumpliendo las siguientes condiciones:
    \begin{enumerate}
        \item $\kappa$ es una fórmula prenexa.
        \item $\kappa$ carece de variables individuales libres.
        \item El prefijo de $\kappa$ comienza con un cuantificador universal y termina con un cuantificador particular.
    \end{enumerate}
\end{definicion}

Entonces con la presente definición podemos deducir el siguiente resultado:

\begin{teorema} \label{thm:teoremaIII}
    Si cada $K$-fórmula es refutable o satisfacible, también lo es cualquier fórmula.
\end{teorema}

\begin{proof}
    Sea $\alpha$ una fórmula que no pertenece a $K$. Sea $\mathfrak{x}$ el conjunto de sus variables libres. Como se puede ver
    directamente, si $\alpha$ es refutable (o satisfacible), se sigue la refutabilidad (o satisfacibilidad) de 
    $\exists \mathfrak{x} \alpha$, e igualmente a la inversa. 

    Sea ahora $\pi \varphi$ la forma normal prenexa de $\exists \mathfrak{x} \alpha$, de tal modo que
    \begin{equation}\label{eq:III-1}
        \exists \mathfrak{x} \alpha \leftrightarrow \pi \varphi
    \end{equation}
    es deducible. Además estipulemos que
    \begin{equation}
        \beta = \forall x \pi \exists y (\varphi \wedge \text{Fx} \vee \lnot \text{Fy})
        \footnote{Las variables x e y no deben aparecer en $\pi$.}
    \end{equation}
    
    Entonces
    \begin{equation}\label{eq:III-2}
        \pi \varphi \leftrightarrow \beta
    \end{equation}
    es deducible (por el \autoref{lem:lema-4} y por la deducibilidad de $\forall x \pi \exists y (\varphi \wedge \text{Fx} \vee \lnot \text{Fy})$).
    $\beta$ pertenece a $K$ y, por tanto, es o satisfacible o refutable. Pero por \eqref{eq:III-1} y \eqref{eq:III-2}
    la satisfacibilidad de $\beta$ implica la de $\exists \mathfrak{x} \alpha$ y consiguientemente también la de $\alpha$,
    y lo mismo se puede aplicar para la refutabilidad. Por tanto, concluimos que también $\alpha$ es o satisfacible o refutable.
\end{proof}

Teniendo en cuenta el \autoref{thm:teoremaIII} anterior, basta para demostrar el \autoref{thm:teoremaII} con probar la siguiente:

\begin{proposicion} \label{prop:satisf-refut}
    Cada $K$-fórmula es satisfacible o refutable.
\end{proposicion}

Para ello definimos previamente el concepto de grado de una $K$-fórmula, y probaremos algunos resultados para poder probar la anterior
proposición.

\begin{definicion}
    Llamaremos grado de una $K$-fórmula
    \footnote{También, en el mismo sentido se le puede llamar "grado de un prefijo".}
    al número de series de cuantificaciones universales de su prefijo, separadas unas de otras
    por cuantificadores existenciales.
\end{definicion}

Probamos primeramente el siguiente resultado.

\begin{teorema}\label{thm:teoremaIV}
    Si cada $K$-fórmula de grado $n$ es o refutable o satisfacible, entonces también lo es cada $K$-fórmula de grado $n+1$.
\end{teorema}

\begin{proof}
    Sea $\pi_1 \alpha$ una $K$-fórmula de grado $n+1$. Sea $\pi_1 = \forall \mathfrak{x} \exists \mathfrak{y} \pi_2$ y 
    $\pi_2 = \forall \mathfrak{u} \exists \mathfrak{v} \pi_3$, donde $\pi_2$ tiene grado $n$ y $\pi_3$ el grado $n-1$. Sea además 
    $\text{F}$ una variable predicativa que no aparezca en $\alpha$. Establezcamos:
    \begin{equation}
        \beta = \forall \mathfrak{x}' \exists \mathfrak{y}' \text{F}  \mathfrak{x}' \mathfrak{y}' \wedge 
        \forall \mathfrak{x} \forall \mathfrak{y} (\text{F}  \mathfrak{x} \mathfrak{y} \rightarrow \pi_2 \alpha)
    \end{equation}
    y
    \begin{equation}
        \gamma = \forall \mathfrak{x}' \forall \mathfrak{x} \forall \mathfrak{y} \forall \mathfrak{u} \exists \mathfrak{y}' \exists \mathfrak{v} 
        \pi_3 (\text{F} \mathfrak{x}' \mathfrak{y}' \wedge (\text{F} \mathfrak{x} \mathfrak{y}) \rightarrow \alpha) \footnote{En esta sentencia
        suponemos que las sucesiones de variables $\mathfrak{x}, \mathfrak{x}', \mathfrak{y}, \mathfrak{y}', \mathfrak{u}, \mathfrak{v}$ son disjuntas
        entre sí.} % esta llave sale en rojo y no se por que, pero no da error
    \end{equation}

    Aplicando ahora dos veces el \autoref{lem:lema-4} junto con el \autoref{lem:lema-6}, obtenemos la deducibilidad de 
    \begin{equation} \label{eq:IV-1}
        \beta \leftrightarrow \gamma
    \end{equation}

    Además, es claro que la fórmula 
    \begin{equation} \label{eq:IV-2}
        \beta \rightarrow \pi_1 \alpha
    \end{equation}
    es válida. Ahora bien, $\gamma$ tiene grado $n$, y por tanto es por hipótesis o satisfacible o refutable. Si $\gamma$ es satisfacible, entonces
    también lo es $\pi_1 \alpha$ (por \eqref{eq:IV-1} y \eqref{eq:IV-2}). Si en caso contrario $\gamma$ es refutable, entonces también lo es 
    $\beta$ (por \eqref{eq:IV-1}), es decir, entonces $\lnot \beta$ es deducible. Sustituyendo ahora F por $\pi_2 \alpha$ en $\lnot \beta$, obtenemos
    que en este caso la siguiente sentencia es deducible:
    \begin{equation}
        \lnot (\forall \mathfrak{x}' \exists \mathfrak{y}' \pi_2 \alpha \wedge 
        \forall \mathfrak{x} \forall \mathfrak{y} (\pi_2 \alpha \rightarrow \pi_2 \alpha))
    \end{equation}

    Se puede observar que, naturalmente, 
    \begin{equation}
        \forall \mathfrak{x} \forall \mathfrak{y} (\pi_2 \alpha \rightarrow \pi_2 \alpha)
    \end{equation}
    es deducible, y por ello también lo es $\lnot \forall \mathfrak{x}' \exists \mathfrak{y}' \pi_2 \alpha$, es decir, en este caso $\pi_2 \alpha$ es refutable.
    Por tanto, de hecho $\pi_2 \alpha$ es o refutable o satisfacible.
\end{proof}

Ahora, para acabar de probar la \autoref{prop:satisf-refut}, sólo necesitamos probar el siguiente resultado:

\begin{teorema} \label{thm:teoremaV}
    Cada $K$-fórmula de primer grado es o satisfacible o refutable.
\end{teorema}

La demostración de este teorema requiere algunas definiciones previas, así como algunos resultados derivados de las definiciones que vamos a establecer, por lo que
se deja la presente demostración para más adelante.

Sea $\forall \mathfrak{x} \exists \mathfrak{y} \alpha (\mathfrak{x} ; \mathfrak{y} )$ $-$abreviado como $\pi\alpha$$-$ una fórmula cualquiera de primer grado.
Mediante $\mathfrak{x}$ representamos un r-tuplo de variables, con $\mathfrak{y}$ un s-tuplo. COnsideremos los r-tuplos sacados de la sucesión $\text{x}_0, \text{x}_1,
\text{x}_2,\dots \text{x}_i \dots$ ordenados por la suma creciente de sus índices en una sucesión:
\begin{equation}
    \mathfrak{x}_1 = (\text{x}_0, \text{x}_0, \dots, \text{x}_0), \quad \mathfrak{x}_2 = (\text{x}_1, \text{x}_0, \dots, \text{x}_0), \quad 
    \mathfrak{x}_3 = (\text{x}_0, \text{x}_1, \text{x}_0, \dots, \text{x}_0) \quad \dots
\end{equation}
y definamos una sucesión $\{\alpha_n\}$ de fórmulas derivadas a partir de $\pi\alpha$ del siguiente modo:
\begin{equation}
    \begin{aligned}
        &\alpha_1 = \alpha(\mathfrak{x}_1 ; \text{x}_1, \text{x}_2, \dots, \text{x}_s) \\
        &\alpha_2 = \alpha(\mathfrak{x}_2 ; \text{x}_{s+1}, \text{x}_{s+2}, \dots, \text{x}_{2s}) \wedge \alpha_1 \\
        &\cdots \\ 
        &\alpha_n = \alpha(\mathfrak{x}_n ; \text{x}_{(n-1)s+1}, \text{x}_{(n-1)s+2}, \dots, \text{x}_{ns}) \wedge \alpha_{n-1}    
    \end{aligned}
\end{equation}

Designemos mediante $\mathfrak{y}_n$ el s-tuplo $\text{x}_{(n-1)s+1}, \text{x}_{(n-1)s+2}, \dots, \text{x}_{ns}$ de tal modo que:
\begin{equation} \label{eq:def-V}
    \alpha_n = \alpha(\mathfrak{x}_n ; \mathfrak{y}_n) \wedge \alpha_{n-1}
\end{equation}

Además, definimos $\pi_n\alpha_n$ estableciendo como sigue:
\begin{equation}
    \pi_n\alpha_n = \exists \text{x}_0 \exists \text{x}_1 \cdots \exists \text{x}_{ns} \alpha_n
\end{equation}

Como fácilmente se comprueba, en $\alpha_n$ aparecen precisamente las variables desde $\text{x}_0$ hasta $\text{x}_{ns}$, que están también ligadas por el prefijo $\pi_n$. 
Además, es evidente que las variables del r-tuplo $\mathfrak{x}_{n+1}$ ya aparecen en $\pi_n$ (y por tanto son distintas de las que aparecen en $\mathfrak{y}_{n+1}$).
Designemos mediante $\pi_n'$ lo que queda de $\pi_n$ cuando suprimimos las variables de r-tuplo $\mathfrak{x}_{n+1}$. Si nos olvidamos del orden de aparición de las variables,
tenemos que $\exists \mathfrak{x}_{n+1} \pi_n' = \pi_n$.

Supuestas todas estas notaciones previas, tenemos como resultado directo el teorema siguiente:

\begin{teorema} \label{thm:teoremaVI}
    Para cada $n$ es deducible $\pi\alpha \rightarrow \pi_n\alpha_n$.
\end{teorema}
\begin{proof}
    Probaremos el teorema mediante inducción.
    \begin{itemize}
        \item Para $n = 1$ tenemos que $\pi\alpha \rightarrow \pi_1\alpha_1$ es deducible, ya que, por el \autoref{lem:lema-3} y la cuarta regla de inferencia, tenemos:
                \begin{equation}
                    \forall \mathfrak{x} \exists \mathfrak{y} \alpha (\mathfrak{x} ; \mathfrak{y} ) \rightarrow 
                    \forall \mathfrak{x}_1 \exists \mathfrak{y}_1 \alpha(\mathfrak{x}_1 ; \mathfrak{y}_1 )
                \end{equation}
                y además, por el \autoref{lem:lema-1} tenemos: 
                \begin{equation}
                    \forall \mathfrak{x}_1 \exists \mathfrak{y}_1 \alpha(\mathfrak{x}_1 ; \mathfrak{y}_1 ) \rightarrow  
                    \exists \mathfrak{x}_1 \exists \mathfrak{y}_1 \alpha(\mathfrak{x}_1 ; \mathfrak{y}_1 )
                \end{equation}
        \item Para un $n$ arbitrario tenemos que $\pi\alpha \wedge \pi_n\alpha_n \rightarrow \pi_{n+1}\alpha_{n+1}$ es deducible ya que, al igual que antes, 
                aplicando el \autoref{lem:lema-3} y la cuarta regla de inferencia, tenemos:
                \begin{equation} \label{eq:VI-1}
                    \forall \mathfrak{x} \exists \mathfrak{y} \alpha (\mathfrak{x} ; \mathfrak{y} ) \rightarrow 
                    \forall \mathfrak{x}_{n+1} \exists \mathfrak{y}_{n+1} \alpha(\mathfrak{x}_{n+1} ; \mathfrak{y}_{n+1} )
                \end{equation}
                y además, por el \autoref{lem:lema-2} tenemos:
                \begin{equation} \label{eq:VI-2}
                    \pi_n\alpha_n \rightarrow \exists \mathfrak{x}_{n+1} \pi_n'\alpha_n
                \end{equation}
                Ahora, aplicamos el \autoref{lem:lema-1} y sustituimos F por $\exists \mathfrak{x}_{n+1} \alpha(\mathfrak{x}_{n+1} ; \mathfrak{y}_{n+1} )$
                y G por $\pi'\alpha_{n}$, con lo que obtenemos:
                \begin{equation} \label{eq:VI-3}
                    \forall \mathfrak{x}_{n+1} \exists \mathfrak{y}_{n+1} \alpha (\mathfrak{x}_{n+1} ; \mathfrak{y}_{n+1} ) \wedge \exists \mathfrak{x}_{n+1} \pi_n'\alpha_n
                    \rightarrow  \exists \mathfrak{x}_{n+1} (\exists \mathfrak{y}_{n+1} \alpha (\mathfrak{x}_{n+1} ; \mathfrak{y}_{n+1} ) \wedge \pi_n'\alpha_n)       
                \end{equation} 

                Fijándonos ahora en que el antecedente del condicional \eqref{eq:VI-3} es la conyunción de los consiguientes de \eqref{eq:VI-1} y \eqref{eq:VI-2},
                obtenemos que es deducible:
                \begin{equation} \label{eq:VI-4}
                    \pi\alpha \wedge \pi_n\alpha_n \rightarrow \exists \mathfrak{x}_{n+1} (\exists \mathfrak{y}_{n+1} 
                    \alpha (\mathfrak{x}_{n+1} ; \mathfrak{y}_{n+1} ) \wedge \pi_n'\alpha_n)
                \end{equation}
                Por otro lado, de \eqref{eq:VI-1} y de los lemas 2, 4 y 6 obtenemos la deducibilidad de:
                \begin{equation} \label{eq:VI-5}
                    \exists \mathfrak{x}_{n+1} (\exists \mathfrak{y}_{n+1} \alpha (\mathfrak{x}_{n+1} ; \mathfrak{y}_{n+1} ) \wedge \pi_n'\alpha_n)
                    \leftrightarrow \pi_{n+1}\alpha_{n+1}
                \end{equation}
                Y gracias a \eqref{eq:VI-4} y \eqref{eq:VI-5} obtenemos la inducción, y por tanto la prueba del teorema.
    \end{itemize}

    Con esto tenemos probado el \autoref{thm:teoremaVI}, que como consecuencia queda probado el \autoref{thm:teoremaV}, que junto con resultados anteriores hemos conseguido 
    probar la \autoref{prop:satisf-refut}. Como vimos anteriormente, esta proposición era el resultado que nos faltaba para acabar la demostración del \autoref{thm:teoremaII}, 
    con lo que hemos probado la tesis de este apartado. Es decir, hemos dado una demostración de que toda fórmula válida de primer orden es deducible.
\end{proof}


\section{Corolarios}




%%%%%%%%%%%%%%%%%%%%%%%%%%%%%%%%%%%%%
\chapter{Sobre sentencias formalmente indecidibles}
\dictum[John von Neumann]{El logro de Gödel en la lógica moderna es singular y monumental - más que monumental, es una señal que 
permanecerá visible lejos en el espacio y en el tiempo.\footnote{Palabras de von Neumann en la entrega a Gödel del premio Einstein 
en 1951, recogidas en \textit{New York Times}, 15 de marzo de 1951, pág. 31.}}


\section{Definiciones y conceptos previos}
Tendremos como objetivo principal probar la existencia de sentencias indecidibles para un sistema formal $P$. Dicho sistema $P$ es
esencialmente el sistema que se obtiene cuando a los axiomas de peano se les añade la lógica de \textit{Principia Mathematica} ($PM$ de aquí en adelante).
\footnote{El hecho de de los axiomas de Peano, así como todas las otras modificaciones del sistema $PM$ introducidas en toda la demostración, sólo 
tienen como finalidad simplificar la prueba, y por ellos son prescindibles.}

Los signos primitivos del sistema $P$ son los siguientes:
\begin{enumerate}
    \item Constantes: « $\sim$ » (no), « $\vee$ » (o), « $\Pi$ » (para todo), « $0$ »(cero), « s » (el siguiente de), «()» (paréntesis).
    \item Variables tipo 1 (para individuos\footnote{Cuando tratamos de individuos hacemos referencia al conjunto de los números naturales.}, 
            incluyendo el 0):  «$\text{x}_1$», «$\text{y}_1$», «$\text{z}_1$», \dots \\
          Variables tipo 2 (para clases de individuos): «$\text{x}_2$», «$\text{y}_2$», «$\text{z}_2$», \dots \\
          Variables tipo 3 (para clases de clases de individuos): «$\text{x}_3$», «$\text{y}_3$», «$\text{z}_3$», \dots \\
          \dots \\
          etc., para cada número natural como tipo.\footnote{Suponemos que disponemos de una cantidad infinita numerable de signos para cada tipo de variables.}
\end{enumerate}
\begin{itemize}
    \item[Observación:] No necesitamos disponer de varaibles para relaciones binarias o $n$-arias ($n>2$) como signos primitivos, ya que podemos definir
                        las relaciones como clase de pares ordenados y los pares ordenados a su vez como clases de clases. Por ejemplo, podemos considerar
                        el par ordenado $<a,b>$ como $\{\{a\}, \{a.b\}\}$, donde $\{x,y\}$ denota la clase cuyos únicos elementos son $x$ e $y$, y $\{x\}$ 
                        la clase cuyo único elemento es $x$.\footnote{Las relaciones no homogéneas también pueden definirse de esta manera; por ejemplo, una
                        relación entre individuos y clases puede definirse como una clase de elementos de la forma $\{\{x_2\}, \{\{x_1\},x_2\}\}$. Todos los 
                        teoremas deducibles en $PM$ son también deducibles cuando se los reformula de esta manera.}
\end{itemize}
                           
Llamaremos \textit{signo de primer tipo} a una combinación de signos que tenga una de las siguientes formas:
\begin{equation}
    a,\thinspace sa,\thinspace ssa,\thinspace sssa,\thinspace \dots,\thinspace \text{etc.,}
\end{equation}
Donde $a$ es $0$ ó es variable de tipo 1. En el primer caso llamamos a tal signo un numeral. Para $n>1$ entendemos por \textit{signo de tipo n} lo mismo que por 
\textit{variable de tipo n}. Llamaremos \textit{fórmulas elementales} a las combianciones de signos de la forma $a(b)$, donde $b$ es un signo de tipo $n$, y $a$ 
es un signo de tipo $n+1$. Definimos la calse de las \textit{fórmulas} como la mínima clase que abarca todas las fórmulas elementales y que, siempre que contenga 
$\alpha$ y $\beta$, contiene también $\sim(\alpha), (\alpha)\vee (\beta)$ y $\Pi\text{x}(\alpha)$ (donde x es una variable cualquiera)
\footnote{Por tanto, $\Pi\text{x}(\alpha)$ es también una fórmula cuando x no aparece o no está libre en $\alpha$. Naturalmente, en ester caso $\Pi\text{x}(\alpha)$
significaría lo mismo que $\alpha$.}. Llamamos a $(\alpha)\vee (\beta)$ la \textit{disyunción} de $\alpha$ y $\beta$, a $\sim(\alpha)$ la \textit{negación} de $\alpha$
y a $\Pi\text{x}(\alpha)$ una \textit{generalización} de $\alpha$. Una \textit{sentencia} es una fórmula sin variables libres (donde la noción de variable libre se 
define del modo usual). A una fórmula con exactamente $n$ variables libres (y ninguna otra variable libre) la llamaremos \textit{signo relacional n-ario}; para $n=1$
lo llamaremos también \textit{signo de clase}.

Por $\daleth^b_v \alpha$ (donde $\alpha$ designa una fórmula, $v$ una variable y $b$ un signo del mismo tipo que $v$) entendemos la fórmula que resulta de reemplazar 
en $\alpha$ cada aparición libre de $v$ por $b$\footnote{Si $v$ no aparece libre en $\alpha$, entonces $\daleth^b_v \alpha = \alpha$. Nótese que $\daleth$ es un signo 
matemático.}. Decimos que una fórmula $\alpha$ es una \textit{elevación de tipo} de otra fórmula $\beta$ si $\alpha$ se obtiene a partir de $\beta$ mediante una elevación
por el mismo número de cada variable que aparece en $\beta$. \\

Las siguientes fórmulas (de I a V) se llaman \textit{axiomas} (están escritas con ayuda de abreviaturas: $\wedge, \supset, \equiv, \Sigma, =$\footnote{Como en $PM$, 
consideramos que $\text{x}_1 = \text{y}_1$ está definido por $\Pi \text{x}_2 (\text{x}_2(\text{x}_1) \supset \text{x}_2(\text{y}_1))$; de igual modo para los tipos 
superiores.}, definidas del modo usual y conforme a las conveciones habituales sobre la omisión de paréntesis)\footnote{Para obtener los axiomas a partir de los esquemas
indicados debemos (después de realizar las sustituciones permitidas en II, III y IV), además,
\begin{enumerate}
    \item[(1)] eliminar las abreviaturas
    \item[(2)] añadir los paréntesis omitidos.
\end{enumerate} 
Nótese que las expresiones así obtenidas deben ser "fórmulas" en el sentido arriba definido.}:

\begin{enumerate}
    \item[I.]   %\begin{equation}
                    %\begin{flalign}
                        $1. \quad \sim (\text{sx}_1 = 0)$ \\
                        $2. \quad \text{sx}_1 = \text{sy}_1 \supset \text{x}_1 = \text{y}_1$ \\
                        $3. \quad \text{x}_2 (0) \wedge \Pi \text{x}_1 \thinspace (\text{x}_2 \thinspace(\text{x}_1) \wedge \text{x}_2 (\text{sx}_1)) \supset \Pi \text{x}_1 \thinspace (\text{x}_2 \thinspace(\text{x}_1)) $   
                    %\end{flalign}
                %\end{equation}
    \item[II.] Cada fórmula que resulta de sustituir X, Y por cualesquiera fórmulas en los siguientes esquemas:
                \begin{flalign}
                    &1.\quad \text{X} \vee \text{X} \supset \text{X} \\
                    &2.\quad \text{X} \supset \text{X} \vee \text{Y} \\
                    &3.\quad \text{X} \vee \text{Y} \supset \text{Y} \vee \text{X} \\
                    &4.\quad (\text{X} \supset \text{Y}) \supset (\text{Z} \vee \text{X} \supset \text{Z} \vee \text{Y})
                \end{flalign}
    \item[III.] Cada fórmula que resulta de uno de estos dos esquemas:
                \begin{flalign}
                    &1.\quad \Pi v\alpha \supset \daleth_v^c \alpha \\ 
                    &2.\quad \Pi v(\beta \vee \alpha) \supset \beta \vee \Pi v(\alpha)
                \end{flalign}
                Cuando sustituimos $\alpha, v, \beta, c$ del siguiente modo (y realizamos la operación indicada por «$\daleth$» en $1.$):
                \begin{enumerate}
                    \item[] Sustituimos por $\alpha$ por una fórmula cualquiera, $v$ por una variable cualquiera, $\beta$ por una fórmula en la que no aparezca libre
                            $v$ y $c$ por un signo del mismo tipo que $v$, siempre que $c$ no contenga alguna variable que pase a estar ligada en un lugar de $\alpha$ 
                            donde $v$ estaba libre.\footnote{Por tanto, $c$ es o una variable o el $0$ o un signo de la forma $s\dots su$, donde $u$ es $0$ o una variable
                            de tipo 1. Respecto de la noción de estar (una variable) libre o ligada en un lugar de $\alpha$, véase \cite{v1927hilbertschen}.}
                \end{enumerate}
    \item[IV.]  Cada fórmula que resulta del esquema $$ \Sigma u \Pi v (u(v) \equiv \alpha) $$ cuando usstituimos $v$ por una varaible cualquiera de tipo $n$, sustituimos
                $u$ por una variable cualquiera de tipo $n+1$ y sustituimos $\alpha$ por una fórmula, en la que $u$ no esté libre. Este axioma desempeña el papel de axioma
                de reducibilidad (el axioma de comprensión de la teoría de conjuntos).
    \item[V.]   Cada fórmula que resulta de $$\Pi \text{x}_1 (\text{x}_2 (\text{x}_1) \equiv \text{y}_2(\text{x}_1)) \supset \text{x}_2 = \text{y}_2$$ por elevación de tipo
                (así como esta fórmula misma). Este axioma dice que una clase está completamente determinada por sus elementos.
\end{enumerate}

Una fórmula $\gamma$ se llama una \textit{inferencia inmediata} de $\alpha$ y $\beta$, si $\alpha$ es la fórmula $\sim \beta \vee \gamma$ (y $\gamma$ se llama una 
\textit{inferencia inmediata} de $\alpha$, si $\gamma$ es la fórmula $\Pi v \alpha $, donde $v$ designa una variable cualquiera). La clase de las \textit{fórmulas deducibles}
se define como la mínima clase de fórmulas que contiene los axiomas y está clausurada respecto a la relación de «inferencia inmediata»\footnote{La regla de sustitución 
resulta aquí supreflua, pues en los axiomas mismo ya tenemos realizadas todas las sustituciones posibles (véase \cite{v1927hilbertschen})}.

Ahora asignamos unívocamente números naturas a los signos primitivos del sistema $P$ del siguiente modo:
\begin{equation}
    \begin{aligned}
        &\text{«}0\text{»} \thinspace \dots \thinspace 1 \\
        &\text{«}s\text{»} \thinspace \dots \thinspace 3 \\
        &\text{«}\sim\text{»} \thinspace \dots \thinspace 5 \\
        &\text{«}\vee\text{»} \thinspace \dots \thinspace 7 \\
        &\text{«}\Pi\text{»} \thinspace \dots \thinspace 9 \\
        &\text{«}(\text{»} \thinspace \dots \thinspace 11 \\
        &\text{«})\text{»} \thinspace \dots \thinspace 13 \\    
    \end{aligned}
\end{equation}

A las variables de tipo $n$ asignamos los números de la forma $\rho^n$ (donde $\rho$ es un número primo $> 13$). Mediante esta asignación a cada fila finita de 
signos primitivos (y en especial a cada fórmula) corresponde biunívocamente una secuencia finita de números naturales. Ahora asignamos (de nuevo biunívocamente)
números naturales a las secuencias finitas de números naturales, haciendo corresponder a la secuencia $n_1, n_2, \dots, n_k$ el número 
$2^{n_1} \cdot 3^{n_2} \cdots \rho^{n_k}_k$ donde $\rho_k$ denota el $k$-avo número primo (en orden de magnitud creciente). Así asignamos biunívocamente un
número natural no sólo a cada signo primitivo, sino también a cada secuencia finita de signos primitivos. Mediante $nu(\text{a})$ denotamos
el número natural asignado al signo primitivo (o a la secuencia de signos primitivos) a. Supongamos que esté dada cierta clase o relación $n$-aria $R$ entre signos
primitivos. Le asignamos la clase o relación $n$-aria $R'$ entre números naturales, en la que están los números $x_1, x_2, \dots, x_n$ si y sólo si hay signos primitivos
o secuencias de signos primitivos $\text{a}_1, \text{a}_2, \dots, \text{a}_n$, tales que $x_i = nu(\text{a}_i)$ (para $i = 1, 2, \dots, n)$ y 
$\text{a}_1, \text{a}_2, \dots, \text{a}_n$ están en la relación $R$. Las clases y relaciones de números naturales, que corresponden de este modo a los conceptos
metamatemáticos hasta ahora definidos, como por ejemplo «variable», «fórmula», «sentencia», «axioma», «fórmula deducible», etc., serán designadas por las mismas 
palabras escritas con letras mayúsculas pequeñas. Por ejemplo, el enunciado de que en el sistema $P$ hay problemas indecidibles se convierte en la siguiente afirmación:
Hay SENTENCIAS $a$, tales que ni $a$ ni la NEGACIÓN de $a$ son FÓRMULAS DEDUCIBLES.\\

Ahora vamos a insertar aquí una disgresión que por el momento no tiene nada que ver con el sistema formal $P$. Empecemos por dar la siguiente definición: 
Decimos que una función numérica\footnote{Es decir, su dominio definicional es la clase de los números naturales (o de los $n$-tuplos de números naturales) 
y sus valores son números naturales.} $f(x_1, x_2, \dots, x_n)$ está \textit{recursivamente definida a partir} de las funciones numéricas $h(x_1, x_2, \dots, x_{n-1})$
y $g(x_1, x_2, \dots, x_{n+1})$, si para cada $x_2, \dots, x_n$, $k$ vale lo siguiente:\footnote{En lo sucesivo las letras latinas minúsculas (a veces con subíndices) 
son siempre variables para números naturales (a no ser que se indique explícitamente lo contrario).}
\begin{equation} \label{eq:recdef}
    \begin{aligned}
        &f(0, x_2, \dots, x_n) = h(x_2, \dots, x_n) \\
        &f(k+1, x_2, \dots, x_n) = g(k, f(k, x_2, \dots, x_n), x_2, \dots, x_n) 
    \end{aligned}
\end{equation}

Una función numérica $f$ se llama \textit{recursiva primitiva} si hay una secuencia finita de funciones $f_1, f_2, \dots, f_n$, que acaba con $f$ y que tiene 
la propiedad de que cada función $f_k$ de la secuencia está recursivamente definida a partir de dos de las funciones precedentes o resulta de alguna de las 
funciones precedentes por sustitución\footnote{Más precisamente: por introducción de algunas de las funciones precedentes en los lugares argumentales de una 
de las funciones precedentes, por ejemplo, $f_k(x_1, x_2) = f_p(f_q(x_1, x_2), f_r(x_2))$, con $p, q, r < k$. No es necesario que todas las variables
del lado izquierdo aparezcan también en el derecho.}, o, finalmente, es una constante o la función del siguiente, $x+1$. La longitud de la mínima secuencia de $f_i$
corresponde a una función recursiva primitiva $f$ se llama su \textit{grado}. Una relación $n$-ara $R$ entre números naturales se llama recursiva 
primitiva\footnote{Incluimos las clases entre las relaciones (como relaciones monarias). Las relaciones recursivas primitivas tienen desde luego la propiedad 
de que para cada $n$-tuplo dado de números naturales $x_1, x_2, \dots, x_n$ \textit{se puede decidir si} $R$ $x_1, x_2, \dots, x_n$ o no.} si hay una función $n$-aria
recursiva primitiva $f$ tal que para cada $n$ números naturales $x_1, x_2, \dots, x_n$ 
\begin{equation}
    R x_1, x_2, \dots, x_n \leftrightarrow f(x_1, x_2, \dots, x_n) = 0 
    \footnote{Para todas las consideraciones intuitivas (y especialmente para las metamatemáticas) usamos el simbolismo de Hilbert definido en \cite{hilbert1962elementos}.}
\end{equation}

Los siguientes teoremas valen:

\begin{teorema} \label{teo:TeoremaI}
    Cada función (o relación) obtenida a partir de funciones (o relaciones) recursivas primitivas por sustitución de las variables por funciones recursivas 
    primitivas es recursiva primitiva; igualmente lo es cada función obtenida a partir de funciones recursivas primitivas por definición recursiva según el
    esquema de \eqref{eq:recdef}.
\end{teorema}

\begin{teorema} \label{teo:TeoremaII}
    Si $R$ y $S$ son relaciones recursivas primitivas, también lo son $\lnot R$ y $R \vee S$ (por tanto, también $R \wedge S$).
\end{teorema}

\begin{teorema} \label{teo:TeoremaIII}
    Si las funciones $f(\mathfrak{x})$, $h(\mathfrak{y})$ son recursivas primitivas, entonces también lo es la relación 
    $f(\mathfrak{x}) = h(\mathfrak{y})$\footnote{Usamos las letras alemanas $\mathfrak{x}, \mathfrak{y}$ como abreviaturas para cualesquiera $n$-tuplos de variables, 
    como por ejemplo $x_1, x_2, \dots, x_n$.}.
\end{teorema}

\begin{teorema} \label{teo:TeoremaIV}
    Si la función $f(\mathfrak{x})$ y la relación $Rx, \mathfrak{y}$ son recursivas primitivas, también lo son las relaciones $S$ y $T$ definidas por 
    \begin{equation}
        \begin{aligned}
            &S(\mathfrak{x}, \mathfrak{y}) \leftrightarrow \exists x (x \leq f(\mathfrak{x}) \wedge Rx, \mathfrak{y} ) \\
            &T(\mathfrak{x}, \mathfrak{y}) \leftrightarrow \forall x (x \leq h(\mathfrak{x}) \wedge Rx, \mathfrak{y} )
        \end{aligned}
    \end{equation}
    así como la función
    \begin{equation}
        q(\mathfrak{x}, \mathfrak{y}) = \mu x (x \leq f(\mathfrak{x}) \wedge Rx, \mathfrak{y}),
    \end{equation}
    donde $\mu x \varphi(x)$ significa el mínimo número $x$, para el que vale $\varphi(x)$, si hay algún tal, y $0$, si no lo hay.
\end{teorema}

El \autoref{teo:TeoremaI} se sigue inmediatamente de la definición de «recursivo primitivo». Tanto el \autoref{teo:TeoremaII} como el \autoref{teo:TeoremaIII}
se basan en que las funciones numéricas
\begin{equation}
    ne(x), \\
    di(x,y), \\
    id(y,x)
\end{equation}
correspondientes a las nociones lógicas $\lnot, \vee, =$, a saber
\begin{flalign}
    &ne(x) = 1; \quad ne(x) = 0 \text{ para } x \neq 0 \\
    &di(0,x) = di(x,0) = 0; \quad di(x,y) = 1 \text{ si } x \neq 0, y \neq 0 \\
    &id(x,y) = 0 \text{ si } x=y; \quad id(x,y) = 1 \text{ si } x \neq y 
\end{flalign}

son recursivas primitivas, como fácilmente se comprueba. He aquí la prueba resumida del \autoref{teo:TeoremaIV}:
\begin{proof}
    Por hipótesis hay una función recursiva primitiva $r(x, \mathfrak{y})$, tal que
    \begin{equation}
        R x, \mathfrak{y} \leftrightarrow r(x, \mathfrak{y}) = 0.
    \end{equation}

    Definamos ahroa mediante el esquema de recursión \eqref{eq:recdef} una función $j(x, \mathfrak{y})$ del siguiente modo:
    \begin{equation}
        \begin{aligned}
            &j(0, \mathfrak{y}) = 0 \\
            &j(n+1, \mathfrak{y}) = (n+1) \cdot a + j(n, \mathfrak{y}) \cdot ne(a) 
            \footnote{Presuponemos como ya es sabido que las funciones $x+y$ (adición) y $x\cdot y$ (multiplicación) son recursivas primitivas.}
        \end{aligned}
    \end{equation}
    donde $a = ne(ne(r(0, \mathfrak{y}))) \cdot ne(r(n+1, \mathfrak{y})) \cdot ne(j(n, \mathfrak{y}))$.

    Por tanto, $j(n+1, \mathfrak{y})$ es igual a $n+1$ (si $a=1$) o es igual a $j(n, \mathfrak{y})$ (si $a=0$)\footnote{$a$ no puede tomar otros valores que
    $0$ y $1$, como se sigue en la definición de $ne$.}. Evidentemente, el primer caso ocurre si y sólo si todos los factores de $a$ son $1$, es decir, si ocurre que
    \begin{equation}
        \lnot R \thinspace 0, \mathfrak{y} \wedge R \thinspace n+1, \mathfrak{y} \wedge j(n, \mathfrak{y}) = 0
    \end{equation}

    De aquí se sigue que la función $j(n, \mathfrak{y})$, considerada como función de $n$, da siempre $0$ hasta (pero no incluyendo) el mínimo valor de $n$
    para el que ocurre $R \thinspace n, \mathfrak{y}$, y a partir de ahí siempre da ese valor. (Por consiguiente, si ya ocurre que $R \thinspace 0, \mathfrak{y},
    j(n, \mathfrak{y})$ es constante e igual a $0$). Por tanto, ocurre que
    \begin{flalign}
        &q(\mathfrak{x}, \mathfrak{y}) = j(f(\mathfrak{x}), \mathfrak{y}) \\
        &S \thinspace \mathfrak{x}, \mathfrak{y} \leftrightarrow R \thinspace q(\mathfrak{x}, \mathfrak{y}), \mathfrak{y} 
    \end{flalign}

    La relación $T$ puede ser reducida, por negación, a un caso análogo al de $S$. Con esto queda probado el \autoref{teo:TeoremaIV}.
\end{proof}

Las funciones $x+y, x\cdot y, x^v$, así como las relaciones $x<y$ y $x=y$ son recursivas primitivas, como fácilmente se comprueba. Partiendo de estos conceptos, 
vamos a definir u n a secuencia de funciones (o relaciones) $1$-$45$, cada una de las cuales se define a partir de las precedentes mediante los procedimientos indicados
en los cuatro teoremas anteriores. En la mayor parte de estas definiciones condensamos en un solo paso varios de los pasos permitidos por estos cuatro teoremas. 
Por tanto, cada una de las funciones (o relaciones) $1$-$45$, entre las que se encuentran, por ejemplo, los conceptos «FÓRMULA», «AXIOMA» e «INFERENCIA INMEDIATA»,
es recursiva primitiva.

\begin{enumerate}
    \item $x/y \leftrightarrow \exists z (z \leq x \wedge x = y \cdot z)$\footnote{Después del definiendum, el signo $=$ se usa en el sentido de «igualdad por
            definición»; el signo $\leftrightarrow$, en el de «equivalencia por definición» (por lo demás, el simbolismo es el de Hilbert).}\\
            $x$ es divisible por $y$.\footnote{Cada vez que en las definiciones siguientes aparece uno de los signos $\forall x, \exists x, \mu x$, éste está 
            seguido de una acotación de $x$. Esta acotación sirve meramente para asegurar que la noción definida es recursiva primitiva (véase \autoref{teo:TeoremaIV}).
            La extensión de la noción definida, por el contrario, no cambiaría en la mayor parte de los casos, aunque dejásemos de lado dicha acotación.}
    \item $\text{Prim } x \leftrightarrow \lnot \exists z (z \leq x \wedge z \neq 1 \wedge z \neq x \wedge x/z) \wedge x>1$\\ $x$ es un número primo.
    \item $0 \thinspace Pr \thinspace x = 0$\\ 
            $(n+1) \thinspace Pr \thinspace x = \mu y (y \leq x \wedge \text{Prim } y \wedge x/y \wedge y>n \thinspace Pr \thinspace x)$ \\
            $n \thinspace Pr \thinspace x$ es el $n$-avo número primo (por orden de magnitud creciente) contenido en $x$\footnote{Para $0 < n \leq z$, donde
            $z$ es el número de diferentes factores primos de $x$. Obsérvese que $n \thinspace Pr \thinspace x = 0$ para $n = z+1$.}.
    \item $0! = 1$ \\ $(n+1)! = (n+1) \cdot n!$
    \item $Pr(0) = 0$ \\ $Pr(n+1) = \mu y (y \leq (Pr(n))! + 1 \wedge \text{Prim } y \wedge y > Pr(n))$ \\ $Pr(n)$ es el $n$-avo número primo (por orden de 
            magnitud creciente)
    \item $n \thinspace Gl \thinspace x = \mu y (y \leq x \wedge x/(n \thinspace Pr \thinspace x)^y \wedge \lnot x/(n \thinspace Pr \thinspace x)^{y+1})$ \\ 
            $n \thinspace Gl \thinspace x$ es el miembro $n$-avo de la secuencia numérica correspondiente al número $x$ (para $n>0$ y $n$ no mayor que la 
            longitud de esa secuencia).
    \item $l(x) = \mu y (y \leq x \wedge y \thinspace Pr \thinspace x > 0 \wedge (y+1) \thinspace Pr \thinspace x = 0)$ \\ $l(x)$ es la longitud de la secuencia 
            numérica correspondiente a $x$.
    \item   $\begin{aligned}
                x * y = \mu z (&z \leq (Pr(l(x) + l(y)))^{x+y} \wedge \forall n (n \leq l(x) \rightarrow n \thinspace Gl \thinspace z = n \thinspace Gl \thinspace x) \wedge \\ 
                &\wedge \forall n (0 < n \leq l(y) \rightarrow (n + l(x)) \thinspace Gl \thinspace z = n \thinspace Gl \thinspace y))
            \end{aligned}$ \\
            $x*y$ corresponde a la operación de concatenación de dos secuencias finitas de números.
    \item $R(x) = 2^x$ \\ $R(x)$ corresponde a la secuencia numérica que sólo consta del número $x$ (para $x>0$).
    \item $E(x) = R(11) * x * R(13)$ \\ $E(x)$ corresponde a la operación de poner entre paréntesis (11 y 13 son los números asignados a los 
            signos primitivos «(» y «)» ).
    \item $n text{ Var } x \leftrightarrow \exists z (13 < z \leq x \wedge text{Prim}(z) \wedge x = z^n) \wedge n \neq 0$ \\ $x$ es una VARIABLE DE TIPO $n$.
    \item $text{ Var } x \leftrightarrow \exists n (n \leq x \wedge n text{ Var } x )$ \\ $x$ es una VARIABLE.
    \item $text{Neg}(x) = R(5) * E(x)$ \\ $text{Neg}(x)$ es la NEGACIÓN de $x$.
    \item $x \text{ Dis } y = E(x) * R(7) * E(y)$ \\ $x \text{ Dis } y$ es la DISYUNCIÓN de $x$ e $y$.
    \item $x \text{ Gen } y = R(9) * R(x) * E(y)$ \\ $x \text{ Gen } y$ es la GENERALIZACIÓN de $y$ respecto a la variable $x$ (suponiendo que $x$ sea una variable).
    \item $0 \thinspace N \thinspace x = x$ \\ $(n+1) \thinspace N \thinspace x = R(3) * n \thinspace N \thinspace x$ \\ 
            $n \thinspace N \thinspace x$ corresponde a la operación de poner $n$ veces el signo $s$ delante de $x$.
    \item $Z(n) = n \thinspace N \thinspace (R(1))$ \\ $Z(n)$ es el NUMERAL que designa el número $n$.
    \item $\text{Typ}_1' \thinspace x \leftrightarrow \exists m \thinspace n(m,n \leq x \wedge (m=1 \vee 1 \text{ Var } m) 
            \wedge x = n \thinspace N \thinspace (R(m)))$\footnote{$m,n \leq x$ es una abreviatura de $m \leq x \wedge n \leq x$ (y lo mismo para 
            más de dos variables)} \\ $x$ es un SIGNO DE TIPO $1$.
    \item $\text{Typ}_n \thinspace x \leftrightarrow (n=1 \wedge \text{Typ}_1' \thinspace (x)) \wedge (n>1 \wedge \exists v (v \leq x \wedge n \text{ Var } v \wedge x = R(v)))$ \\
            $x$ es un SIGNO DE TIPO $n$.
    \item $Elf \thinspace x \leftrightarrow \exists yzn (y, z, n \leq x \wedge \text{Typ}_n(y) \wedge \text{Typ}_{n+1}(z) \wedge x = x* E(y))$ \\ $x$ es una FÓRMULA ELEMENTAL.
    \item $Op \thinspace x \thinspace y \thinspace z \leftrightarrow x = \text{Neg}(y) \wedge x = y \text{ Dis } z \vee \exists v (v \leq x \wedge \text{Var } v \wedge 
            x = v \text{ Gen } y)$ 
    \item $\begin{aligned}
            FR \thinspace x \leftrightarrow \forall n (0 < n \leq l(x) \rightarrow Elf(n \thinspace Gl \thinspace x) &\vee \exists pq (0 < p,q < n \wedge \\
            \wedge &Op(n \thinspace Gl \thinspace x, p \thinspace Gl \thinspace x,q \thinspace Gl \thinspace x))) \wedge l(x) > 0
        \end{aligned}$ \\
            $x$ es una SECUENCIA DE FÓRMULAS, cada una de las cuales es o una FÓRMULA ELEMENTAL o se obtiene de las precedentes mediante las operaciones de 
            NEGACIÓN, DISYUNCIÓN o GENERALIZACIÓN.
    \item $\text{Form } x \leftrightarrow \exists n (n \leq (Pr(l(x)^2))^{x \cdot (l(x)^2)} \wedge FR \thinspace n \wedge x = (l(x)) \thinspace Gl \thinspace n)$\footnote{La
            acotación $n \leq (Pr(l(x)^2))^{x \cdot (l(x)^2)}$ puede comprobarse así: La longitud de la más corta secuencia de  fórmulas correspondientes a $x$ puede ser a lo 
            sumo igual al número de subfórmulas de $x$. Pero hay a lo sumo $l(x)$ subfórmulas de longitud $1$, a lo sumo $l(x) -1$ de longitud $2$, \dots, por tanto en conjunto
            a lo sumo $\frac{l(x) \cdot (l(x) + 1)}{2} \leq n \leq (l(x))^2$. Por tanto podemos suponer que todos los factores primos de $n$ son menores que $Pr((l(x))^2)$, 
            que su número es $\leq (l(x))^2$ y que sus exponentes (que son subfórmulas de $x$) son $\leq x$.} \\ 
            $x$ es una FÓRMULA (es decir, el último miembro de una SECUENCIA DE FÓRMULAS $n$).
    \item $v \thinspace \text{Geb} \thinspace n, x \leftrightarrow \text{Var } v \wedge \text{Form } x \wedge \exists a \thinspace b \thinspace c \thinspace 
            (a, b, c \leq x \wedge x = a*(v \text{ Gen } b)* c \wedge \text{Form } x \wedge l(a) + 1 \leq n \leq l(a) + l(v \text{ Gen } b))$ \\ La VARIABLE $v$
            está LIGADA en $x$ en el $n$-avo lugar.
    \item $v \thinspace Fr \thinspace n, x \leftrightarrow \text{Var } v \wedge \text{Form } x \wedge v = n \thinspace Gl \thinspace x \wedge n \leq l(x) 
            \wedge \lnot v \text{ Geb } n, x$ \\ La VARIABLE $v$ está LIBRE en $x$  en el $n$-avo lugar.
    \item $v \thinspace Fr \thinspace x \leftrightarrow \exists n (n \leq l(x) \wedge v \thinspace Fr \thinspace n, x)$ \\ $v$ aparece como VARIABLE LIBRE EN $x$.
    \item $Su \thinspace x\tbinom{n}{y} = \mu z (z \leq (Pr(l(x) + l(y)))^{x+y} \wedge \exists u \thinspace v(u,v \leq x \wedge x = u * R(n \thinspace Gl \thinspace x) * v \wedge \\ 
            \wedge z =  u * y * v \wedge n = l(u) + 1))$ \\
            $Su \thinspace x\tbinom{n}{y}$ se obtiene a partir de $x$ cuando sustituimos el $n$-avo miembro de $x$ por $y$ (suponiendo que $0 < n \leq (x)$).
    \item $\qquad \quad 0 \thinspace St \thinspace v, x = \mu n (n \leq l(x) \wedge v \thinspace Fr \thinspace n, x \wedge \lnot \exists p(n < p \leq l(x) \wedge v \thinspace Fr \thinspace p, x))$\\
            $(k+1) \thinspace St \thinspace v, x = \mu n (n < k \thinspace St \thinspace v, x \wedge v \thinspace Fr \thinspace n, x \wedge \lnot \exists 
            p(n < p \leq k \thinspace St \thinspace v, x \wedge v \thinspace Fr \thinspace p, x))$\\
            $k \thinspace St \thinspace v, x$ es el $k+1$-avo lugar de $x$ (contado a partir del extremo derecho de la FÓRMULA $x$), en el que $v$ aparece LIBRE en $x$ 
            (y es $0$ si no hay tal lugar).
    \item $A(v,x) = \mu n (n \leq l(x) \wedge n \thinspace St \thinspace v, x = 0)$ \\ $A(v,x)$ es el número de lugares en que $v$ aparece LIBRE en $x$.
    \item $Sb_0 (x^v_y) = x$ \\ $Sb_{k+1} (x^v_y) = Su \thinspace (Sb_k (x^v_y)) \tbinom{\text{\textit{k St v,x}}}{y}$
    \item $Sb (x^v_y) = Sb_{A(v,x)} (x^v_y)$\footnote{Si $v$ no es una VARIABLE  o $x$ no es una FÓRMULA, enotnces $Sb (x^v_y) = x$.}\\
            $Sb (x^v_y)$ es la noción anteriormente definida de $\daleth^b_v \alpha$\footnote{En vez de $Sb(Sb(x^v_w)^w_z)$ escribimos $Sb \thinspace x^{v \thinspace w}_{y \thinspace z}$, 
            (y análogamente para más de dos VARIABLES).}.
    \item $x \text{ Imp } y = (\text{Neg}(x)) \text{ Dis } y \\ x \text{ Con } y = \text{Neg } ((\text{Neg}(x)) \text{ Dis } (\text{Neg}(y))) \\
            x \text{ Aeq } y = (x \text{ Imp } y) \text{ Cong}(y \text{ Imp } x) \\ v \text{ Ex } y = \text{Neg } (v \text{ Gen} (\text{Neg}(y))).$
    \item $ n \thinspace Th \thinspace x = \mu y (y \leq x^{(x^n)} \wedge \forall k (k \leq l(x) \rightarrow (k \thinspace Gl \thinspace x > 13 
            \wedge k \thinspace Gl \thinspace y = k \thinspace Gl \thinspace x \cdot (1 \thinspace Pr(k \thinspace Gl \thinspace x))^n)))$ \\
            $ n \thinspace Th \thinspace x $ es la $n$-ava ELEVACIÓN DE TIPO de $x$ (suponiendo que $x$ y $ n \thinspace Th \thinspace x $ sean fórmulas).
            \begin{itemize}
                \item[Nota. ] A los axiomas I,$1$-$3$, corresponden tres números determinados, que designamos mediante $z_1$, $z_2$ y $z_3$. 
                                Procedemos a definir los siguientes:
            \end{itemize}
    \item $Z - Ax \thinspace x \leftrightarrow x = z_1 \vee x = z_2 \vee x = z_3$
    \item $A_1 - Ax \thinspace \leftrightarrow \exists y (y \leq x \wedge \text{Form } y \wedge x = (y \text{ Dis } y) \text{ Imp } y)$ \\
            $x$ es una FÓRMULA que se obtiene a partir del esquema axiomático II,$1$ por sustitución. De modo análogo se definen $A_2-Ax$, 
            $A_3-Ax$ y $A_4-Ax$, correspondientes a los axiomas II,$2$-$4$. 
    \item $A-Ax \thinspace x \leftrightarrow A_1-Ax \thinspace x \vee A_2-Ax \thinspace x \vee A_3-Ax \thinspace x \vee A_4-Ax \thinspace x$\\
            $x$ es una FÓRMULA que se obtiene por sustitución a partir de un esquema axiomático conectivo. 
    \item $Q \thinspace z, y, v \leftrightarrow \lnot \exists n \thinspace m \thinspace w \thinspace (n \leq l(y) \wedge m \leq l(z) \wedge 
            w \leq z \wedge w = m \thinspace Gl \thinspace z \wedge w \text{ Geb } n,y \wedge v \thinspace Fr \thinspace n, y)$ \\
            $z$ no contiene VARIABLE alguna que esté LIGADA en $y$ en un lugar, en el cual $v$ esté LIBRE.
    \item $L_1 - Ax \thinspace x \leftrightarrow \exists v\thinspace y \thinspace z \thinspace n (v,y,z,n \leq x \wedge n \text{ Var } v 
            \wedge \text{Typ}_n z \wedge \text{Form } y \wedge Q \thinspace z,y,v \wedge \\ \wedge x = (v \text{ Gen } y) \text{ Imp } (Sb(y_z^v)))$\\ 
            $x$ es una FÓRMULA que se obtiene por sustitución a partir del esquema axiomático III,$1$.
    \item $L_2 - Ax \thinspace x \leftrightarrow \exists v\thinspace q \thinspace p \thinspace (v,p,q \leq x \wedge n \text{ Var } v 
            \wedge \text{Form } p \wedge \lnot v \thinspace Fr \thinspace p \wedge \text{Form } q \wedge \\
            \wedge x = (v \text{ Gen } (p \text{ Dis } q)) \text{ Imp } (p \text{ Dis } (v \text{ Gen } q)))$\\ 
            $x$ es una FÓRMULA que se obtiene por sustitución a partir del esquema axiomático III,$2$.
    \item $R-Ax \thinspace x \leftrightarrow \exists u \thinspace v \thinspace y \thinspace n \thinspace (u,v,y,n \leq x \wedge 
            n \text{ Var } v \wedge (n+1) \text{ Var } u \wedge \lnot u \thinspace Fr \thinspace y \wedge \text{Form } y \wedge 
            x = u \text{ Ex} (v \text{ Gen } ((R(u) * E(R(v))) \text{ Aeq } y)))$ \\ 
            $x$ es una FÓRMULA que se obtiene por sustitución a partir del esquema axiomático IV,$1$.
            \begin{itemize}
                \item[Nota. ] Al axioma V,$1$ le corresponde un número determinado $z_4$. Así, procedemos a definir:
            \end{itemize}
    \item $M-Ax \thinspace x \leftrightarrow \exists n(n\leq x \wedge x = n \thinspace Th \thinspace z_4)$
    \item $Ax \thinspace x \leftrightarrow Z-Ax \thinspace x \vee L_1-Ax \thinspace x \vee L_2-Ax \thinspace x \vee R-Ax \thinspace x \vee M-Ax \thinspace x$\\ 
            $x$ es un AXIOMA.
    \item $Fl \thinspace x \thinspace y \thinspace z \leftrightarrow y = z \text{ Imp } x \vee \exists v (v \leq x 
            \wedge \text{Var } v \wedge x = v \text{ Gen } y)$\\ 
            $x$ es una INFERENCIA INMEDIATA de $y$ y $z$.
    \item $Bw \thinspace x \leftrightarrow \forall n (0 < n \leq l(x) \rightarrow Ax(n \thinspace Gl \thinspace x) \wedge \exists pq (0 < p,q < n \wedge 
            Fl \thinspace (n \thinspace Gl \thinspace x, p \thinspace Gl \thinspace x, q \thinspace Gl \thinspace x))) \wedge l(x) > 0 $\\
            $x$ es una DEDUCCIÓN (una secuencia finita de FÓRMULAS, cada una de las cuales es o un AXIOMA o una INFERENCIA INMEDIATA 
            de dos de las FÓRMULAS precedentes).
    \item $x \thinspace B \thinspace y \leftrightarrow Bw \thinspace x \wedge (l(x)) \thinspace Gl \thinspace x = y$ \\ $x$ es una DEDUCCIÓN DE LA FÓRMULA y.
    \item $\text{Bew } x \leftrightarrow \exists y \thinspace yBx$ \\ $x$ es una FÓRMULA DEDUCIBLE. ($\text{Bew } x$ es la única de las nociones 
            $1$-$46$ de la que no podemos afirmar que sea recursiva primitiva).
\end{enumerate}

El hecho puede ser formulado vagamente diciendo que cada relación recursiva primitiva es definible en el sistema $P$ (interpretado en cuanto a su 
contenido del modo habitual) puede ser expresado con precisión y sin referencia a ninguna interpretación natural de las fórmulas de $P$, 
mediante el siguiente teorema:

\begin{teorema} \label{teo:TeoremaV}
    Para cada relación recursiva primitiva $n$-aria $R$ hay un SIGNO RELACIONAL $r$ (con las VARIABLES LIBRES\footnote{Las VARIABLES $u_1, \dots, u_n$ pueden
    estar elegidas de cualquier manera.} $u_1, u_2, \dots, u_n$ tal que para cada $n$-tuplo de números naturales $(x_1, \dots, x_n)$ vale:
    
    \begin{equation} \label{eq:Teo5-1}
        R x_1, \dots, x_n \rightarrow \text{Bew} \left(Sb \left(r 
        \begin{aligned}
            &u_1, \dots, u_n\\
            &Z(x_1), \dots, Z(x_n)
        \end{aligned}\right)\right)
    \end{equation}

    \begin{equation} \label{eq:Teo5-2}
        \lnot R x_1, \dots, x_n \rightarrow \text{Bew} \left(\text{Neg } Sb \left(r 
        \begin{aligned}
            &u_1, \dots, u_n\\
            &Z(x_1), \dots, Z(x_n)
        \end{aligned}\right)\right)
    \end{equation}
\end{teorema}

\begin{proof}
    Aquí nos limitaremos a esbozar la prueba de este teorema\footnote{Desde luego, el \autoref{teo:TeoremaV} se basa en el hecho de que, dada una relación
    recursiva primitiva $R$, para cada $n$-tuplo de números naturales podemos decidir en base a los axiomas del sistema $P$ si esos números están en la 
    relación $R$ o no.}, pues no ofrece ninguna dificultad de principio, pero es bastante larga. Probamos el teorema para todas las relaciones
    $R x_1, \dots, x_n $ de la forma $x_1 = f(x_2, \dots, x_n)$\footnote{De ahí se sigue inmediatamente su validez para toda relación recursiva
    primitiva, pues cada tal relación es equivalente a $0 = f(x_1, \dots, n_n)$, donde $f$ es recursiva primitiva.} (donde $f$ es una función recursiva 
    primitiva) por inducción completa sobre el grado de $f$. Para funciones de grado $1$ (es decir, constantes y la función $x + 1$) el teorema es trivial. 
    Sea $f$ de grado $m$. Entonces $f$ se obtiene a partir de las funciones $f_1, \dots, f_k$ de menor grado mediante las operaciones de sustitución o de 
    definición recursiva. Puesto que el teorema ya está probado para $f_1, \dots, f_k$, por hipótesis inductiva, hay SIGNOS RELACIONALES correspondientes
    $r_1, \dots, r_k$, tales que \eqref{eq:Teo5-1} y \eqref{eq:Teo5-2} valen. Los procesos de definición (sustitución y definición recursiva) mediante los 
    que se obtiene $f$ a partir de $f_1, \dots, f_k$ pueden ser ambos formalmente deducidos en el sistema $P$. Si se hace esto se obtiene a partir de $r_1, \dots, r_k$
    un nuevo SIGNO RELACIONAL $r$\footnote{Desde luego, cuando esta prueba se lleva a cabo con todo detalle, $r$ no se define indirectamente con ayuda 
    de su interpretación intuitiva, sino sólo por su estructura puramente formal.} para el cual puede probarse sin dificultad que valen \eqref{eq:Teo5-1} 
    y \eqref{eq:Teo5-2}, haciendo uso de la hipótesis inductiva. Un SIGNO RELACIONAL $r$, que corresponde\footnote{Que, por tanto, en la interpretación 
    usual expresa que esa relación se da.} de este modo a una relación recursiva primitiva, se llama recursivo primitivo.
\end{proof}

Ahora llegamos a la meta de nuestras consideraciones. Sea $K$ una clase cualquiera de FÓRMULAS. Designamos mediante $\text{Flg}(K)$\footnote{conjunto de 
inferencias a partir de $K$.} el mínimo conjunto de FÓRMULAS que contiene todas las FORMULAS de $K$ y todos los AXIOMAS y está clausurado respecto a 
la relación de «INFERENCIA INMEDIATA». Decimos que $K$ es $\omega$-consistente si no hay ningún SIGNO DE CLASE $a$, tal que
\begin{equation}
    \forall n \left(
        Sb \left( a
            \begin{aligned}
                v \\
                Z(n)
            \end{aligned}
        \right)
        \in \text{Flg}(K)
    \right)
    \wedge (\text{Neg}(v \text{ Gen } a)) \in \text{Flg}(K)
\end{equation}
donde $v$ es la VARIABLE LIBRE del SIGNO DE CLASE $a$. 

Cada sistema $\omega$-consistente es también consistente, desde luego. Pero la inversa no vale, como más adelante veremos.

El resultado general sobre la existencia de sentencias indecidibles dice como sigue:

\begin{teorema} \label{teo:TeoremaVI}
    Para cada clase recursiva primitiva y $\omega$-consistente $K$ de FORMULAS hay un SIGNO DE CLASE $r$ tal que ni $v \text{ Gen } a$ ni 
    $\text{Neg}(v \text{ Gen } a)$ pertenecen a $\text{Flg}(K)$ (donde $v$ es la VARIABLE LIBRE de $r$).
\end{teorema}

\begin{proof}
    
\end{proof}
\endinput
%------------------------------------------------------------------------------------
% FIN DEL CAPÍTULO. 
%------------------------------------------------------------------------------------



% --------------------------------------------------------------------
% APPENDIX: Opcional
% --------------------------------------------------------------------

\appendix % Reinicia la numeración de los capítulos y usa letras para numerarlos
\pdfbookmark[-1]{Apéndices}{appendix} % Alternativamente podemos agrupar los apéndices con un nuevo \part{Apéndices}

% !TeX root = ../libro.tex
% !TeX encoding = utf8

\chapter{Primer apéndice}\label{ap:apendice1}

Los apéndices son opcionales.

Archivo: \texttt{apendices/apendice01.tex}

\endinput
%------------------------------------------------------------------------------------
% FIN DEL APÉNDICE. 
%------------------------------------------------------------------------------------

% Añadir tantos apéndices como sea necesario 

% --------------------------------------------------------------------
% GLOSARIO: Opcional
% --------------------------------------------------------------------

% !TeX root = ../libro.tex
% !TeX encoding = utf8

% \chapter*{Glosario}
% \addcontentsline{toc}{chapter}{Glosario} % Añade el glosario a la tabla de contenidos

% La inclusión de un glosario es opcional.

% Archivo: \texttt{glosario.tex}

% \begin{description} 
%   \item[$\mathbb{R}$] Conjunto de números reales.

%   \item[$\mathbb{C}$] Conjunto de números complejos.

%   \item[$\mathbb{Z}$] Conjunto de números enteros.
% \end{description}
% \endinput
 

% -------------------------------------------------------------------
% BACKMATTER
% -------------------------------------------------------------------

\backmatter % Desactiva la numeración de los capítulos
\pdfbookmark[-1]{Referencias e Índices}{BM-Referencias}

% BIBLIOGRAFÍA
%-------------------------------------------------------------------

\setbibpreamble{Las referencias se listan por orden alfabético. Aquellas referencias con más de un autor están ordenadas de acuerdo con el primer autor.\par\bigskip}
\bibliographystyle{alpha} 
\begin{small} % Normalmente la bibliografía se imprime en un tamaño de letra más pequeño.
\bibliography{library.bib}
\end{small}


% ÍNDICE TERMINOLÓGICO  (Opcional) 
%------------------------------------------------------------------- 

% Para incluir el índice terminológico es necesario descomentar los siguientes comandos. Incluir un índice terminológico es opcional

% \cleardoublepage 
% \begin{footnotesize} % Normalmente el índice se imprime en un tamaño de letra más pequeño.
% \printindex 
% \end{footnotesize}

\end{document}
